\documentclass[11pt,a4paper]{article}
\usepackage{fontspec}
\usepackage{polyglossia}
\setmainlanguage{greek}
\setotherlanguage{english}
\defaultfontfeatures{Ligatures=TeX}
\setmainfont{FreeSerif}
\newfontfamily\greekfont[Ligatures=TeX]{FreeSerif}
\setsansfont{FreeSerif}
\newfontfamily\greekfontsf[Ligatures=TeX]{FreeSerif}
\setmonofont{DejaVu Sans Mono}
\newfontfamily\greekfonttt{DejaVu Sans Mono}
\usepackage{amsmath,amssymb}
\usepackage[a4paper,top=2.5cm,bottom=2.5cm,left=2.6cm,right=2.6cm]{geometry}
\usepackage{titlesec}
\titleformat{\section}{\normalfont\large\bfseries}{\thesection.}{0.5em}{}
\titleformat{\subsection}{\normalfont\normalsize\bfseries}{\thesubsection}{0.5em}{}
\titlespacing*{\section}{0pt}{16pt}{6pt}
\titlespacing*{\subsection}{0pt}{10pt}{4pt}
\usepackage{fancyhdr}
\pagestyle{fancy}\fancyhf{}
\renewcommand{\headrulewidth}{0pt}
\fancyfoot[C]{\small\thepage}
\usepackage[hidelinks]{hyperref}
\usepackage{enumitem}
\setlist{itemsep=2pt,topsep=4pt,leftmargin=1.6em}
\setlength{\parindent}{1.2em}
\setlength{\parskip}{2pt}
\linespread{1.1}\setlength{\emergencystretch}{3em}\Urlmuskip=0mu plus 1mu
\newcommand{\Rs}{R^{\star}}
\newcommand{\Mth}{\mathcal{M}_{\text{θέσις}}}

\begin{document}

\begin{center}
{\small ΕΛΛΗΝΙΚΟ ΜΕΣΟΓΕΙΑΚΟ ΠΑΝΕΠΙΣΤΗΜΙΟ · Τμήμα Μουσικής Τεχνολογίας και Ακουστικής}\par\vspace{0.8cm}
{\large\bfseries ΜΟΥΣΙΚΟΛΟΓΙΚΟ ΥΠΟΜΝΗΜΑ}\par\vspace{0.4cm}
{\Large\bfseries Η μουσική θέση ως υπολογιστική μονάδα\\[2pt] της βυζαντινής μουσικής}\par\vspace{0.3cm}
{\large Μουσικολογικό πλαίσιο, μεθοδολογικές αρχές\\ και παλαιογραφικές προεκτάσεις του ερευνητικού σχεδίου}\par\vspace{0.6cm}
\begin{tabular}{@{}ll@{}}
Υποψήφιος & Γεώργιος Σπυριδάκης\\
Συνοδεύει & Σχέδιο διδακτορικής έρευνας (βασικός πυρήνας και επεκτάσεις)\\
Τεκμηρίωση & \href{https://github.com/GeorgeSpy/byz-position-phd-study}{github.com/GeorgeSpy/byz-position-phd-study}\\
Ημερομηνία & Οκτώβριος 2026\\
\end{tabular}
\end{center}
\vspace{0.4cm}

%========================================================== 1
\section{Το μουσικολογικό πρόβλημα: η θέση και η ταυτότητά της μέσα στην παραλλαγή}

Η θεωρία της βυζαντινής ψαλτικής περιγράφει ήδη από τον 15ο αιώνα το μέλος όχι ως απλή παράθεση μεμονωμένων φθόγγων, αλλά ως σύνθεση καθιερωμένων μελικών σχηματισμών, των θέσεων. Ο Μανουήλ Χρυσάφης ο λαμπαδάριος εντάσσει τη θέση στην τέχνη της μελοποιίας (έκδ. Conomos 1985, σσ.~40--41), ενώ ο Γαβριήλ ιερομόναχος εισάγει την εύγλωττη αναλογία της λέξης: όπως τα γράμματα συντίθενται σε λέξεις και οι λέξεις σε λόγο, έτσι και τα επιμέρους φωνητικά και άφωνα σημάδια συντίθενται σε θέσεις που συγκροτούν το ψαλλόμενο μέλος (έκδ. Hannick και Wolfram 1985, σ.~72). Στην παλαιότερη, συνοπτικότερη σημειογραφία, η θέση δεν ταυτιζόταν με μια σταθερή διαδοχή διαστηματικών βημάτων. Όπως συνοψίζει ο Ε. Γιαννόπουλος (2014, σ.~161), τα φωνητικά σημάδια έδειχναν την ποσοτική κίνηση της φωνής πάνω στις βαθμίδες, ενώ το ζητούμενο μέλος προέκυπτε από τη σύμπραξή τους με τα άφωνα σημάδια, τις μεγάλες υποστάσεις (πρβλ. Αλεξάνδρου 2010· Alexandru 2012· Αλεξάνδρου 2017). Ο συγκεκριμένος συνδυασμός διαστηματικών και χειρονομικών σημαδιών πάνω σε ορισμένη βαθμίδα της κλίμακας γεννούσε τη μελωδική φράση, τη θέση, ανάλογα με το γένος, το γένος μελοποιίας και τον δρόμο του μέλους, ενώ η αναλυτική του εκφορά μεταδιδόταν προφορικά και ακουστικά μέσω της διδασκαλίας.

Στο ίδιο συμπέρασμα καταλήγει, από τη σκοπιά της μορφολογικής και παλαιογραφικής ανάλυσης του παπαδικού ρεπερτορίου, ο N. Gheorghiță (2009· 2010), εξετάζοντας συστηματικά σαράντα δύο κυριακά Κοινωνικά του 18ου αιώνα, των Δανιήλ Πρωτοψάλτη, Πέτρου Πελοποννησίου και Πέτρου Βυζαντίου, στην παλαιότερη γραφή και στις αναλυτικές εξηγήσεις τους στη Νέα Μέθοδο. Ορίζοντας ως θεμελιώδη μονάδα του μέλους τη φόρμουλα, δηλαδή τη σύζευξη μιας θέσης με το χειρονομικό σημάδι που την κυβερνά, απομονώνει έντεκα χαρακτηριστικούς τύπους του παπαδικού ιδιώματος: κύλισμα, παρακάλεσμα, παρακλητική, επέγερμα, αργοσύνθετον--γοργοσύνθετον, ξηρόν κλάσμα, πίασμα, σύναγμα, τρομικοπαρακάλεσμα και δύο σύνθετους συνδυασμούς σημαδιών (Gheorghiță 2009, σ.~127). Η σύγκριση των πηγών δείχνει ότι η μελική μορφή της ίδιας θέσης μεταβάλλεται συστηματικά υπό την επίδραση τεσσάρων παραγόντων (ό.π., σσ.~250--251): (α) της βαθμίδας μέσα στο τετράχορδο ή το πεντάχορδο του ήχου, καθώς η θέση μετατίθεται με τον τροχό ή με μετάθεση, γεγονός που είχε ήδη επισημάνει ο Χρύσανθος (1832, §408), σημειώνοντας ότι το παρακάλεσμα γεννά διαφορετικό μέλος στον Πα και στον Βου· (β) του γένους, διατονικού, χρωματικού ή εναρμονίου, το οποίο μεταμορφώνει λ.χ. την καταληκτική θέση με κύλισμα· (γ) του γένους μελοποιίας, ειρμολογικού, στιχηραρικού ή παπαδικού, δεδομένου ότι η ιδιαίτερη μελική ενέργεια ορισμένων θέσεων ατονεί έξω από το παπαδικό μέλος· και (δ) του χρώματος του χειρονομικού σημαδιού, ερυθρού ή μελανού, στους παλαιούς κώδικες, το οποίο, σύμφωνα και με τις θεωρητικές πραγματείες του Αποστόλου Κώνστα του Χίου (Αποστολόπουλος 2002), υποδείκνυε αν η θέση έπρεπε να αναπτυχθεί απλή ή σύνθετη, συντομότερη ή εκτενέστερη.

Όταν το παπαδικό μέλος καταγράφεται αναλυτικά στη Νέα Μέθοδο του 1814, αυτή η εγγενής πλαστικότητα της θέσης δημιουργεί μια κρίσιμη ασυμμετρία ανάμεσα στην επιφανειακή ομοιότητα της γραφής και στη μουσικολογική ταυτότητα του σχηματισμού. Από τη μία πλευρά, χωρία με την ίδια ή πολύ συγγενή μετροφωνία, δηλαδή τη διαστηματική ανάγνωση των ποσοτικών σημείων, μπορεί να αντιστοιχούν σε διαφορετικές ιστορικές θέσεις. Στη μελέτη του για την παράδοση και την εξήγηση των Χερουβικών, ο Κ.~Χ. Καραγκούνης (2003, σ.~788) τεκμηριώνει περιπτώσεις όπου η ίδια μετροφωνική επιφάνεια εξηγεί διαφορετικούς σχηματισμούς της παλαιότερης γραφής, οι οποίοι διακρίνονταν εκεί μέσω μεγάλων υποστάσεων όπως το αντικένωμα, το λύγισμα και το έτερον παρακάλεσμα. Από την άλλη πλευρά, η ίδια θέση μπορεί να εμφανίζεται με αισθητά διαφοροποιημένη αναλυτική επιφάνεια, είτε επειδή προσαρμόζεται στην ενεργή μαρτυρία ώστε να διατηρηθεί σταθερό το τονικό σημείο εκκίνησης ή κατάληξης, όπως παρατηρείται λ.χ. στην εναρκτήρια θέση του Χερουβικού του Μανουήλ Αγαλλιανού (Καραγκούνης 2005, σ.~175), είτε επειδή μεταβάλλεται ο αριθμός των συλλαβών, ο λεκτικός τονισμός και η αναλυτικότητα της καταγραφής (Αρβανίτης 2010, τ.~Α΄, §5.2, σσ.~136--142· Καραγκούνης 2003, σσ.~810--812· πρβλ. Σπανουδάκης 2017 για τα κεκραγάρια). Επιπλέον, τα όρια μιας θέσης μέσα στη συνεχή ροή του παπαδικού μέλους δεν συμπίπτουν μηχανικά με τα όρια των συλλαβών ή των λέξεων: μια εναρκτήρια θέση μπορεί να εκτείνεται σε τρεις διαδοχικές συλλαβές («Οἱ τὰ Χε-»), ενώ σε εκτενή μελίσματα μία μόνο συλλαβή φιλοξενεί περισσότερες διαδοχικές θέσεις και βραχείς συνδετικούς σχηματισμούς (Καραγκούνης 2003, σσ.~192, 207, 209).

Η διαπίστωση αυτή έχει άμεση συνέπεια τόσο για τη μουσικολογία όσο και για την υπολογιστική ανάλυση μουσικής, την ανάκτηση μουσικής πληροφορίας (Music Information Retrieval, MIR). Ένας αλγόριθμος που αναζητά απλώς επαναλαμβανόμενες ακολουθίες διαστημάτων ή μετρά την επιφανειακή ομοιότητα δύο αποσπασμάτων θα αποτύχει ακριβώς στις μουσικολογικά ενδιαφέρουσες περιπτώσεις: θα συγχέει ομοιομετροφωνικές αλλά διαφορετικές θέσεις και θα αδυνατεί να συνδέσει τις σύντομες, τις αναπτυγμένες ή τις τροπικά μετατεθειμένες πραγματώσεις της ίδιας θέσης. Αντίστροφα, η βυζαντινή ψαλτική προσφέρει στην υπολογιστική έρευνα ένα προνομιακό πεδίο που σπάνια απαντά σε άλλες μουσικές παραδόσεις, διότι διαθέτει δύο ιστορικά τεκμηριωμένους μηχανισμούς μετασχηματισμού του ίδιου μουσικού αντικειμένου: αφενός τη μελοποιητική επανάληψη και προσαρμογή της θέσης από έργο σε έργο, κατά τους παράγοντες που ανέδειξε η μορφολογική έρευνα, και αφετέρου την εξήγηση, δηλαδή την αναλυτική απόδοση του ίδιου παλαιότερου χωρίου στη Νέα Μέθοδο, συχνά μάλιστα από δύο σύγχρονους εξηγητές, τον Γρηγόριο Πρωτοψάλτη και τον Χουρμούζιο Χαρτοφύλακα. Η κεντρική ιδέα της προτεινόμενης διατριβής είναι ότι η ταυτότητα μιας μουσικής θέσης μπορεί να μελετηθεί υπολογιστικά ως ό,τι παραμένει αναγνωρίσιμο μέσα από αυτούς τους δύο τεκμηριωμένους μετασχηματισμούς, υπό την προϋπόθεση ότι η μουσικολογική γνώση θα ορίσει ρητά τη μονάδα ανάλυσης, τις διαγνωστικές της σχέσεις και τα ανεξάρτητα κριτήρια αξιολόγησης.

%========================================================== 2
\section{Εννοιολογικές διακρίσεις και επίπεδα τεκμηρίωσης}

Για να μετατραπεί η παραπάνω μουσικολογική προβληματική σε αυστηρό ερευνητικό πρωτόκολλο, η πρόταση διαχωρίζει εννοιολογικά τέσσερις όρους που στην καθημερινή αναλυτική πρακτική συχνά συνεκφέρονται: τη θέση, την εμφάνιση, την πραγμάτωση και την οικογένεια θέσεων.

Ως μουσική θέση νοείται ο ιστορικά και θεωρητικά αναγνωρίσιμος μελικός σχηματισμός μεταβλητής έκτασης. Ως εμφάνιση ορίζεται η εντοπισμένη παρουσία μιας θέσης σε συγκεκριμένο χωρίο ενός έργου και ενός ιστορικού μάρτυρα, χειρογράφου ή εντύπου, με ρητά σημεία αρχής και τέλους· επειδή στη μελισματική ροή η ακριβής τομή μεταξύ δύο διαδοχικών θέσεων ενδέχεται να επιδέχεται περισσότερες από μία μουσικολογικά νόμιμες αναγνώσεις, η εμφάνιση δεν περιορίζεται αναγκαστικά σε ένα μοναδικό ζεύγος ορίων, αλλά μπορεί να καταγράφεται με σύνολο αποδεκτών εναλλακτικών ορίων. Ως πραγμάτωση ορίζεται η συγκεκριμένη μελωδική, χρονική, τροπική και προσωδιακή μορφή που λαμβάνει η θέση στη δεδομένη εμφάνιση: το μήκος της σε μελωδικά γεγονότα, η κατανομή της στις συλλαβές, οι εσωτερικοί καλλωπισμοί ή οι συντμήσεις της. Ο όρος «πραγμάτωση» αφορά εδώ τη γραπτή μελική μορφή στη σημειογραφία της πηγής και διακρίνεται ρητά από τη φωνητική εκτέλεση του χωρίου από έναν ψάλτη σε ηχογράφηση.

Εμφανίσεις από διαφορετικά έργα ή πηγές οι οποίες κρίνονται μουσικολογικά ως πραγματώσεις του ίδιου μελικού τύπου συγκροτούν μια οικογένεια θέσεων. Η διάκριση ανάμεσα στην ιστορική ονομασία μιας θέσης και στην αναλυτική κατηγορία της οικογένειας είναι μεθοδολογικά αναγκαία: στις πηγές και στη διδασκαλική παράδοση, η ίδια ιστορική ονομασία μπορεί να καλύπτει σχηματισμούς που έχουν διαφοροποιηθεί σε διακριτούς μελικούς τύπους, ενώ αντίστροφα οργανικά συγγενείς πραγματώσεις μπορεί να απαντούν χωρίς σταθερή ονομασία. Για τον λόγο αυτό, κάθε οικογένεια θέσεων στη διατριβή δεν ορίζεται από την επιφανειακή ομοιότητα ούτε απλώς από ένα όνομα, αλλά από μια γραπτή, αριθμημένη προδιαγραφή οικογένειας. Η προδιαγραφή κωδικοποιεί σε ελέγξιμη μορφή τα κριτήρια που χρησιμοποιεί ο μουσικολόγος, με αφετηρία την καταγωγή από μεγάλο σημάδι, τη λειτουργική θέση μέσα στη φράση, αρκτική, συνδετική ή καταληκτική, και την τροπική συμπεριφορά:
\begin{enumerate}
\item τις αναγκαίες σχέσεις, οι οποίες πρέπει να ικανοποιούνται σε μια ενεργή διαμόρφωση του τύπου (λ.χ. κατάληξη σε συγκεκριμένη βαθμίδα του ενεργού ήχου με ορισμένη καταληκτική προσέγγιση)·
\item τις υποστηρικτικές σχέσεις και τις εναλλακτικές επαρκείς διαμορφώσεις, οι οποίες επιτρέπουν στην ίδια οικογένεια να περιλαμβάνει λ.χ. τόσο μια σύντομη όσο και μια αναπτυγμένη ή χρωματικά μετατεθειμένη μορφή, χωρίς να απαιτείται ένας μοναδικός, άκαμπτος πυρήνας πανομοιότυπων φθόγγων (πολυθετική δομή)· και
\item τις αποκλειστικές σχέσεις, δηλαδή ρητά κριτήρια, όπως μια συγκεκριμένη φθορά ή μετατόπιση της βάσης, των οποίων η παρουσία αποκλείει την υπαγωγή στην οικογένεια ακόμη και αν η μετροφωνική επιφάνεια είναι όμοια.
\end{enumerate}

Η ανάγκη να διακρίνεται η ιστορική ονομασία ενός χειρονομικού σημαδιού από την αναλυτική κατηγορία της οικογένειας θέσεων αναδεικνύεται ανάγλυφα στην αντιπαραβολή του κυλίσματος και της παρακλητικής από τον Gheorghiță (2009, §7.2, σσ.~130--133, 207--209). Στα εξεταζόμενα Κοινωνικά του 18ου αιώνα, το κύλισμα εμφανίζεται κατά κανόνα στις καταληκτικές περιοχές του ψαλμικού στίχου, όπου ένα συνοπτικό διαστηματικό περίγραμμα της παλαιότερης γραφής αναπτύσσεται μέσω της εξήγησης σε εκτενές μέλισμα, ενώ σε μη καταληκτικές θέσεις η μελισματική του ενέργεια περιορίζεται αισθητά. Αντίθετα, η παρακλητική συμμετέχει σε ποικίλους νευματικούς συνδυασμούς που επιτελούν διαφορετικές μορφολογικές λειτουργίες: απαντά κυρίως σε εναρκτήριες φράσεις, αλλά και σε περιοχές κρατημάτων, ενώ η ανάπτυξή της διαφοροποιείται μεταξύ παπαδικού και ειρμολογικού ιδιώματος. Συνεπώς, για την υπολογιστική και μουσικολογική ανάλυση, η παρουσία του ίδιου χειρονομικού σημαδιού στην παλαιά γραφή, ή της ίδιας ονομασίας στην αναλυτική παράδοση, παρέχει μεν ισχυρή υποψήφια συγγένεια, αλλά δεν προδικάζει αν όλες οι εμφανίσεις της παρακλητικής συγκροτούν μία ενιαία οικογένεια με εναλλακτικές λειτουργικές διαμορφώσεις ή αν διασπώνται σε διακριτές οικογένειες. Το ερώτημα αυτό δεν λύνεται αξιωματικά, αλλά μετατρέπεται σε αντικείμενο εμπειρικού ελέγχου στις μελέτες Μ1, Α2 και Μ3a, παρέχοντας παράλληλα ιδανικό ιστορικό υλικό για τη συγκρότηση «δύσκολων θετικών» και «δύσκολων αρνητικών» ζευγών αξιολόγησης.

Η διατριβή δεν προϋποθέτει εκ των προτέρων ότι οι οικογένειες αυτές αποτελούν κλειστές, αντικειμενικά δεδομένες οντότητες. Αντίθετα, θέτει την ίδια τη συνοχή τους υπό εμπειρικό έλεγχο: αν ανεξάρτητοι ειδικοί της ψαλτικής, κρίνοντας τυφλά νέα χωρία, συμφωνούν τόσο στην ονομασμένη θέση όσο και στην οικογενειακή υπαγωγή βάσει της προδιαγραφής, τεκμηριώνεται ότι η κατηγορία αποτυπώνει πραγματική διακριτική γνώση της παράδοσης· αν συμφωνούν στην ευρύτερη θέση αλλά αποκλίνουν συστηματικά στην επιμέρους οικογένεια, η οικογένεια καταγράφεται ως αναλυτική κατασκευή με ρητά όρια εγκυρότητας. Και στις δύο περιπτώσεις, το αποτέλεσμα συνιστά καθαρό μουσικολογικό εύρημα.

Τέλος, στο ιστορικό επίπεδο, η εξήγηση, δηλαδή η αναλυτικότερη απόδοση παλαιότερου μέλους στη Νέα Μέθοδο ή σε ενδιάμεσες εξηγητικές γραφές του 18ου και των αρχών του 19ου αιώνα, εξετάζεται παράλληλα με τις συγγενείς τεχνικές του καλλωπισμού και της σύντμησης (Gheorghiță 2009, κεφ.~4.1· Στάθης 1978). Όταν για το ίδιο παλαιότερο χωρίο σώζονται δύο ή περισσότερες αναλυτικές αποδόσεις, λ.χ. του Γρηγορίου και του Χουρμουζίου, κάνουμε λόγο για συν-εξήγηση. Η ιδιότητα δύο αποδόσεων ως συν-εξηγήσεων δεν συνάγεται ποτέ από το πόσο μοιάζουν μελωδικά, αλλά τεκμηριώνεται αποκλειστικά με φιλολογική και πηγολογική κριτική: ταύτιση του εξηγούμενου έργου, του λειτουργικού κειμένου και της θέσης της συλλαβής μέσα στη σύνθεση, και έλεγχος της γενεαλογικής σχέσης ή τυχόν αντιγραφικής εξάρτησης των κωδίκων.

Για να αποτραπεί κάθε σύγχυση ανάμεσα στα δεδομένα των πηγών, στις υποθέσεις του ερευνητή και στις μετρήσεις των υπολογιστικών μοντέλων, ολόκληρη η βάση δεδομένων της διατριβής οργανώνεται σε τρία διαχωρισμένα επίπεδα τεκμηρίωσης. Το επίπεδο της πηγής περιλαμβάνει αποκλειστικά ό,τι παραδίδει ο ιστορικός μάρτυρας: τα γραφικά σημάδια της σημειογραφίας, τις μαρτυρίες και τις φθορές, το ποιητικό κείμενο, τις επιγραφές (αποδιδόμενο μελοποιό, εξηγητή, χρονολογία) και τις τυχόν ρητές ονομασίες θέσεων ή μεγάλων υποστάσεων που σημειώνει ο γραφέας. Το επίπεδο της ανάλυσης περιλαμβάνει τις μουσικολογικές κρίσεις: τα προτεινόμενα όρια μιας εμφάνισης, την ένταξή της σε οικογένεια βάσει συγκεκριμένης έκδοσης προδιαγραφής, τον χαρακτηρισμό ενός χωρίου ως αμφίσημου ή ανεπίλυτου και την πηγολογική σύνδεση δύο χωρίων ως συν-εξηγήσεων· κάθε εγγραφή φέρει το αναγνωριστικό του αναλυτή και τον ρόλο της στη μελέτη, ως υλικό ορισμού προδιαγραφής, υλικό ανάπτυξης ή υλικό ανεξάρτητης τυφλής αξιολόγησης. Το επίπεδο της παρατήρησης περιλαμβάνει εξωτερικές εμπειρικές μετρήσεις που παράγονται ανεξάρτητα από τη μουσικολογική ανάλυση της παρτιτούρας, όπως οι μετρήσεις χρονισμού και τονικού ύψους από ηχογραφήσεις σύγχρονων ψαλτών στα ίδια χωρία ή οι αποκρίσεις ακροατών σε ψυχοακουστικές δοκιμασίες. Ο διαχωρισμός αυτός εξασφαλίζει ότι μια αναλυτική υπόθεση δεν παρουσιάζεται ποτέ ως πρωτογενές δεδομένο του κώδικα ούτε χρησιμοποιείται κυκλικά για να επιβεβαιώσει τον εαυτό της.

%========================================================== 3
\section{Η κατάσταση των πηγών και η αρχιτεκτονική του ερευνητικού σχεδίου}

Μια κρίσιμη μεθοδολογική απόφαση της πρότασης είναι η διάκριση ανάμεσα στον βασικό ερευνητικό πυρήνα, ο οποίος είναι αυτάρκης και υλοποιήσιμος εντός των 48 μηνών, και στις επεκτάσεις υπό προϋποθέσεις, οι οποίες ενεργοποιούνται κλιμακωτά μόνο εφόσον η απογραφή των πηγών τεκμηριώσει επαρκή όγκο και ανεξαρτησία δεδομένων. Η διάκριση αυτή επιβάλλεται από την άνιση διαθεσιμότητα των μουσικολογικών πηγών στο παπαδικό ρεπερτόριο.

Για τη μελέτη των θέσεων σε πλήρη παπαδικά έργα της Νέας Μεθόδου, Χερουβικά και Κοινωνικά, υπάρχει πλούσιο και άμεσα προσβάσιμο υλικό. Ως καταλογογραφική αφετηρία και ευρετήριο προέλευσης αξιοποιείται η ψηφιακή συλλογή \emph{Papadic Formulae} της Ιεράς Μονής Αγίου Αντωνίου Αριζόνας, η οποία αποδελτιώνει μελωδικούς τύπους από 758 συνθέσεις (349 Χερουβικά, 351 Κοινωνικά και 58 λοιπά μέλη), οργανωμένους ανά ήχο. Επειδή ο κατάλογος αυτός αποτυπώνει την πρακτική κατηγοριοποίηση ψαλτών χωρίς ρητό θεωρητικό πρωτόκολλο οριοθέτησης, στη διατριβή δεν χρησιμοποιείται ως αλήθεια αναφοράς, αλλά ως δεξαμενή υποψήφιων σχέσεων και οδηγός εντοπισμού χωρίων. Η κύρια ανάλυση γίνεται πάνω σε πλήρη έργα κωδικοποιημένα ψηφιακά, μέσω του ανοικτού λογισμικού Neanes, σε διαρκή αντιπαραβολή με τις δημοσιευμένες μορφολογικές αναλύσεις (Gheorghiță 2009, 2010· Καραγκούνης 2003, 2005· Ιωαννίδης 2022). Συνεπώς, όλα τα ερωτήματα που αφορούν την κωδικοποίηση, την κοινή οριοθέτηση και αναγνώριση θέσεων (Μελέτη Μ1), τον έλεγχο της αξίας των μουσικολογικών προδιαγραφών (Μελέτη Α2), τη σχέση οικογενειακής ταυτότητας και πραγμάτωσης από εμφανίσεις διαφορετικών έργων (Μελέτη Μ2a), την τυφλή επικύρωση από ειδικούς (Μελέτη Μ3a) και τη βασική χαρτογράφηση πλήρων έργων (Μελέτη Μ5-min) στηρίζονται σε επαρκές υλικό και συγκροτούν τον βασικό πυρήνα.

Αντίθετα, όταν περνούμε σε ερωτήματα που απαιτούν είτε παράλληλες εξηγήσεις του ίδιου χωρίου από δύο εξηγητές είτε άμεση αντιπαραβολή παλαιάς γραφής και αναλυτικής εξήγησης σε επίπεδο μεμονωμένης θέσης, τα δεδομένα στο παπαδικό γένος είναι περιορισμένα. Όπως επισημαίνει ο Γιαννόπουλος (2014, σσ.~178--179), οι σωζόμενες χειρόγραφες συλλογές εξηγήσεων θέσεων αφορούν συνηθέστερα το στιχηραρικό ρεπερτόριο, ενώ για παπαδικές θέσεις οι μαρτυρίες είναι ελάχιστες· στον κώδικα 35 της Μονής Αγίας Αναστασίας Φαρμακολυτρίας, των ετών περίπου 1815--1819, παρατίθενται δώδεκα θέσεις από Χερουβικό β΄ ήχου, από τις οποίες μόνο δύο φέρουν διπλή εξήγηση και τρεις μένουν ανεξήγητες (ό.π., σ.~174). Ο ίδιος σημειώνει ότι σχεδόν δεν σώζονται μεταγραφές μελών του Γρηγορίου και του Χουρμουζίου σε παράλληλη παλαιά και νέα γραφή (σ.~178). Παράλληλα, υπάρχουν συγκεκριμένοι, καλά τεκμηριωμένοι πυρήνες ολοκληρωμένων συνθέσεων: ο Gheorghiță (2009, σσ.~83, 128--129) τεκμηριώνει ότι η σειρά των οκτώ κυριακών Κοινωνικών του Δανιήλ Πρωτοψάλτη εξηγήθηκε από τον Πέτρο Πελοποννήσιο και πέρασε στη Νέα Μέθοδο μέσω του Γρηγορίου Πρωτοψάλτη και του Χουρμουζίου Χαρτοφύλακα, με επιπλέον εξηγήσεις από τον Ιωάσαφ Ριλιώτη, τον Απόστολο Κώνστα και ανώνυμο εξηγητή, σε κώδικες όπως οι BAR gr.~955, Βατοπαιδίου 1289 και 1458, BN Βουκουρεστίου gr.~4920 και στις έντυπες εκδόσεις του \emph{Ταμείου Ανθολογίας}.

Επειδή ο ακριβής αριθμός των ιστορικά ανεξάρτητων ζευγών συν-εξηγήσεων, δηλαδή ζευγών όπου δεν μεσολαβεί απλή αντιγραφή του ενός εξηγητή από τον άλλο, και η τομή τους με διαθέσιμες ηχογραφήσεις πολλών ψαλτών στα ίδια χωρία (όπως, ενδεχομένως, το σώμα ΔΑΜΑΣΚΗΝΟΣ· Delviniotis και Kouroupetroglou 2011) μπορούν να προσδιοριστούν μόνο μετά από συστηματική πηγολογική απογραφή, η πρόταση εντάσσει την απογραφή και τεκμηρίωση των συν-εξηγήσεων στο υποχρεωτικό πρώτο εξάμηνο του προγράμματος (Πακέτο Εργασίας ΠΕ0, μήνες 1--6). Αν η απογραφή αναδείξει επαρκή αριθμό επιλέξιμων ζευγών και ηχογραφήσεων, το ιστορικό σκέλος της εκπαίδευσης (Μ2b: σύγκριση B2--EISE) και οι παλαιογραφικές επεκτάσεις (αντίστροφη ανάκτηση παλαιού σημαδιού ATL-INV και μορφολογία μετασχηματισμών ATL) ενεργοποιούνται σε πλήρη ποσοτική κλίμακα· αν το υλικό είναι αριθμητικά περιορισμένο, εκτελούνται ως τεκμηριωμένες μουσικολογικές μελέτες περίπτωσης, χωρίς να διακυβεύεται η στατιστική και μεθοδολογική αρτιότητα του βασικού πυρήνα.

%========================================================== 4
\section{Συμβολική κωδικοποίηση της Νέας Μεθόδου και διαλειτουργικότητα με το πρότυπο MEI}

Η είσοδος κάθε υπολογιστικού μοντέλου στη συμβολική ανάκτηση μουσικής πληροφορίας προϋποθέτει μια κρίσιμη απόφαση: σε ποια κλίμακα και με ποια γνωρίσματα θα αναπαρασταθεί η μουσική γραφή. Στη δυτική σημειογραφία, η στοιχειώδης μονάδα εισόδου ταυτίζεται σχεδόν αυτονόητα με τον μεμονωμένο φθόγγο, ο οποίος φέρει συγκεκριμένο τονικό ύψος και διάρκεια. Στη σχεσιακή νευματική σημειογραφία της Νέας Μεθόδου, ωστόσο, η σχέση ανάμεσα στο γραφικό σύμβολο και στη μελωδική κίνηση είναι σύνθετη: ένα μεμονωμένο γραφικό σημάδι μπορεί να δηλώνει μία μόνο διαστηματική κίνηση, ενώ ένας συνδυασμός σημαδιών, λ.χ. ένα σύμπλεγμα ποσοτικών, χρονικών και ποιοτικών χαρακτήρων, μπορεί να λειτουργεί ως ενιαία νευματική ομάδα που αποδίδει περισσότερα διαδοχικά μελωδικά γεγονότα. Αν η ψηφιακή κωδικοποίηση κρατήσει μόνο τα διαστηματικά βήματα, χάνει τον τρόπο με τον οποίο ο γραφέας ή ο εξηγητής ομαδοποίησε τα σημάδια· αν κρατήσει μόνο τα ακατέργαστα γραφικά σύμβολα, δυσχεραίνει τη σύγκριση χωρίων όπου η ίδια μελωδική κίνηση έχει γραφτεί με εναλλακτική ορθογραφία.

Για να μην προδικαστεί αυθαίρετα ποια κλίμακα ανάγνωσης είναι η καταλληλότερη για την αναγνώριση των θέσεων, η διατριβή κατασκευάζει ένα ενιαίο συμβολικό επίπεδο αναφοράς (Πακέτο Εργασίας ΠΕ1), πάνω στο οποίο συγκρίνονται συστηματικά τρεις μονάδες εισόδου, $U=\{u_1,u_2,u_3\}$:
\begin{enumerate}
\item το γραφικό σημείο $u_1$, όπως ακριβώς εμφανίζεται στην πηγή·
\item το ατομικό μελωδικό γεγονός $u_2$, δηλαδή η στοιχειώδης διαστηματική και χρονική κίνηση που προκύπτει από την ανάγνωση ενός ή περισσότερων σημείων· και
\item η σύνθετη νευματική μονάδα $u_3=f_{\text{group}}(N)$, η οποία διατηρεί την ομαδοποίηση των νευμάτων μέσω ενός ντετερμινιστικού και αντιστρέψιμου κανόνα που υπολογίζεται αποκλειστικά από τα σημειογραφικά στοιχεία, χωρίς να εισάγει πληροφορία από τις ετικέτες των θέσεων.
\end{enumerate}

Κάθε μία από τις τρεις αυτές μονάδες εισόδου διασταυρώνεται με τέσσερα σωρευτικά σύνολα γνωρισμάτων, $\Phi=\{N,\;N{+}T,\;N{+}T{+}M,\;N{+}T{+}M{+}L\}$, δημιουργώντας ένα πλέγμα δώδεκα κωδικοποιήσεων ($3\times4$). Το βασικό σημειογραφικό στρώμα $N$ περιλαμβάνει τη σχετική διαστηματική κίνηση και την πλήρη ταυτότητα τόσο των ποσοτικών όσο και των ποιοτικών χαρακτήρων, ώστε να μη χάνεται η πληροφορία σημείων όπως το ψηφιστόν, το αντικένωμα ή το ομαλόν. Το χρονικό στρώμα $T$ προσθέτει τους χρονικούς χαρακτήρες (κλάσμα, απλή, γοργόν, δίγοργον κ.ο.κ.) και τις παράγωγες σχετικές χρονικές αξίες τους σε κοινή μονάδα ανάγνωσης, διακρίνοντάς τες ρητά τόσο από τις εκδοτικές διαστολές όσο και από την απόλυτη αγωγή της φωνητικής εκτέλεσης. Το τροπικό στρώμα $M$ ενσωματώνει τον ενεργό ήχο, τις ενδιάμεσες μαρτυρίες και τις φθορές, επιτρέποντας να παρακολουθείται η σχέση της θέσης με τις βαθμίδες και τις μεταβολές του ήχου ή του γένους (Χρύσανθος 1832· Troelsgård 2011· Alexandru και Troelsgård 2020). Τέλος, το κειμενικό και προσωδιακό στρώμα $L$ καταγράφει τη συλλαβική μονάδα, τη λέξη, τον λεκτικό τόνο, τα όρια λέξης και φράσης και την τοπική μελισματική πυκνότητα, δηλαδή τον αριθμό μελωδικών γεγονότων ανά συλλαβή, διακρίνοντας παράλληλα τα λεξικά χωρία από τα μη λεξικά, όπως τα κρατήματα (πρβλ. Αναστασίου 2005). Η σταδιακή αυτή προσθήκη των στρωμάτων $T$, $M$ και $L$ έχει άμεσο μουσικολογικό κίνητρο: επιτρέπει να μετρηθεί πειραματικά σε ποιον βαθμό η αναγνώριση μιας θέσης εξαρτάται από τη σχέση των σημαδοφώνων με τα γειτονικά τους στοιχεία, την τοποθέτησή τους στη φράση και την προσαρμογή τους στον λεκτικό τονισμό (Δεσπότης 2007· Καραγκούνης 2003· Αρβανίτης 2010).

Θεμελιώδης αρχή αυτής της υποδομής είναι η αμφίδρομη ιχνηλασιμότητα και ο έλεγχος αντίστροφης ανακατασκευής. Κάθε μελωδικό γεγονός ή σύνθετη μονάδα διατηρεί αμφίδρομο σύνδεσμο προς τα πρωτογενή γραφικά σημάδια της πηγής, αρχικά μέσω ντετερμινιστικών προσαρμογέων από δομημένα αρχεία του λογισμικού Neanes, ώστε οποιαδήποτε πρόβλεψη ή σφάλμα του αλγορίθμου να ελέγχεται άμεσα πάνω στη σημειογραφία του κώδικα ή της έκδοσης. Πριν χρησιμοποιηθεί οποιαδήποτε κωδικοποίηση, εφαρμόζεται δοκιμή αντίστροφης ανακατασκευής, η οποία επαληθεύει ότι από τη συμβολική ακολουθία ανακτάται χωρίς απώλειες η αρχική σημειογραφική πληροφορία.

Η αρχιτεκτονική αυτή ευθυγραμμίζεται με τις αρχές που έχουν τεθεί διεθνώς για την ψηφιακή κωδικοποίηση της βυζαντινής σημειογραφίας στο πλαίσιο του προτύπου Music Encoding Initiative (MEI). Όπως έχουν αναδείξει οι M.~Alexandru, E.~De~Luca και I.~Fujinaga (2024· πρβλ. Siklafidis και Alexandru 2021) για τη μεσοβυζαντινή γραφή και οι M.~Thomae, M.~Alexandru κ.ά. (2026) για την κωδικοποίηση των νευμάτων της Νέας Μεθόδου (χρυσάνθινης σημειογραφίας) σε MEI, μια έγκυρη αναπαράσταση οφείλει να διατηρεί διακριτά τόσο τη γραφική μορφή και τα ελληνικά ονόματα των σημαδιών όσο και τη διαστηματική, χρονική και ποιοτική τους σημασία, χωρίς να συγχέει το σημειογραφικό δεδομένο με την αναλυτική ερμηνεία. Το ενιαίο συμβολικό επίπεδο της διατριβής σχεδιάζεται ώστε να είναι διαλειτουργικό και συμβατό με το αναπτυσσόμενο σχήμα MEI για τα βυζαντινά νεύματα, προσφέροντας στην ερευνητική κοινότητα ένα μουσικολογικά ελεγμένο σώμα παπαδικών μελών με επισημειωμένες θέσεις, οικογένειες και παράλληλες εξηγήσεις.

%========================================================== 5
\section{Μελέτη Μ1: Η από κοινού οριοθέτηση και αναγνώριση θέσεων και η σύγκριση με τη στατιστική κατάτμηση}

Πάνω στο πλέγμα των δώδεκα κωδικοποιήσεων, η πρώτη μελέτη του βασικού πυρήνα απαντά στο θεμελιώδες ερώτημα: ποια ανάγνωση της γραφής επιτρέπει να εντοπίζουμε ταυτόχρονα πού αρχίζει και πού τελειώνει μια θέση μέσα στη συνεχή ροή του μέλους και σε ποια οικογένεια ανήκει, και σε ποιον βαθμό η μουσικολογική οριοθέτηση των θέσεων συμπίπτει ή αποκλίνει από μια καθαρά στατιστική ή επιφανειακή κατάτμηση;

Στην υπολογιστική βιβλιογραφία, η μελωδική κατάτμηση και η κατηγοριοποίηση τμημάτων αντιμετωπίζονται συνήθως ως δύο διαδοχικά και ανεξάρτητα προβλήματα: πρώτα ένας αλγόριθμος κόβει τη μελωδία σε τμήματα και κατόπιν ένας δεύτερος προσπαθεί να τα ταξινομήσει. Στο παπαδικό μέλος, όμως, τα δύο αυτά ερωτήματα είναι αλληλένδετα. Όταν μια εναρκτήρια θέση εκτείνεται σε τρεις συλλαβές («Οἱ τὰ Χε-»), μια παύση, μια αλλαγή συλλαβής ή ένα τοπικό άλμα θα οδηγούσε έναν απλό κατατμητή να κόψει τη θέση στη μέση· αντίστροφα, όταν μέσα σε ένα εκτενές μέλισμα μιας μόνο συλλαβής διαδέχονται η μία την άλλη δύο ή τρεις διαφορετικές θέσεις (Καραγκούνης 2003, σσ.~207, 209), δεν υπάρχει κανένα κειμενικό όριο για να υποδείξει την τομή, και το σημείο όπου τελειώνει η πρώτη θέση και αρχίζει η δεύτερη γίνεται αντιληπτό μόνο επειδή ο αναλυτής αναγνωρίζει την εσωτερική μορφολογία και των δύο σχηματισμών.

Για να αποτυπώσει υπολογιστικά αυτή τη μουσικολογική πραγματικότητα, η Μ1 ορίζει ως βασικό εκτιμητή ένα ημι-μαρκοβιανό τυχαίο πεδίο υπό συνθήκη (semi-Markov Conditional Random Field, semi-CRF· Sarawagi και Cohen 2004· Kong κ.ά. 2016). Σε αντίθεση με τα συνήθη μοντέλα που αποδίδουν ετικέτα σε κάθε μεμονωμένο φθόγγο χωριστά, το semi-CRF εξετάζει και βαθμολογεί ολόκληρα υποψήφια τμήματα μεταβλητού μήκους $[s,e]$ πάνω στη συνεχή ακολουθία, αξιολογώντας ταυτόχρονα τα χαρακτηριστικά των ορίων τους και τη συνολική εσωτερική τους δομή ως προς κάθε γνωστή οικογένεια θέσεων. Για να ελεγχθεί αν η γνώση της οικογενειακής ταυτότητας βοηθά πράγματι στον εντοπισμό των ορίων, η Μ1 συγκρίνει, με τον ίδιο κωδικοποιητή και προϋπολογισμό, τρεις εσωτερικές αρχιτεκτονικές: την αρχιτεκτονική A (κοινό semi-CRF), όπου η οριοθέτηση και η αναγνώριση της οικογένειας γίνονται ταυτόχρονα και αλληλεπιδρούν κατά την εκπαίδευση· την αρχιτεκτονική B (CRF ορίων και ταξινομητής), όπου πρώτα ανιχνεύονται τα όρια σε επίπεδο γεγονότος χωρίς γνώση της οικογένειας και κατόπιν ταξινομούνται τα τμήματα· και την αρχιτεκτονική C (τμηματική κατάτμηση χωρίς οικογενειακή ανάδραση), όπου ο κατατμητής βαθμολογεί ολόκληρα τμήματα ως προς το αν αποτελούν «θέση γενικώς», χωρίς η οικογενειακή ταξινόμηση να επηρεάζει την εκπαίδευση των ορίων. Η σύγκριση της A με τις B και C δίνει άμεση, ποσοτική απάντηση στο μουσικολογικό ερώτημα αν τα όρια μιας θέσης στη Νέα Μέθοδο είναι απλώς τοπικά σημεία άρθρωσης της μελωδίας ή αν προσδιορίζονται από την ίδια την ταυτότητα του μελικού σχηματισμού.

Ο σχεδιασμός του semi-CRF στη Μ1 ενσωματώνει τέσσερις δικλίδες που σέβονται την ειδική φύση των μουσικολογικών δεδομένων:
\begin{enumerate}
\item \emph{Εκπαίδευση από μερικές επισημειώσεις} (Tsuboi κ.ά. 2008). Σε ένα πλήρες Χερουβικό ή Κοινωνικό, ο αναλυτής μπορεί να επισημειώσει με βεβαιότητα τις εμφανίσεις των εξεταζόμενων οικογενειών, αλλά ανάμεσά τους υπάρχουν συνδετικά χωρία, σπάνιοι σχηματισμοί ή περιοχές που δεν ανήκουν στο αρχικό λεξιλόγιο. Αν το μοντέλο εκπαιδευόταν θεωρώντας κάθε μη επισημειωμένο φθόγγο ως αρνητικό παράδειγμα, θα μάθαινε εσφαλμένα να απορρίπτει πραγματικές θέσεις που απλώς δεν έχουν ακόμη καταγραφεί. Γι' αυτό, η πιθανοφάνεια του semi-CRF αθροίζει πάνω σε όλες τις συμβατές αναλύσεις στις μη επισημειωμένες περιοχές, κρατώντας τες ως λανθάνουσες μεταβλητές.
\item \emph{Σύνολο αποδεκτών εναλλακτικών ορίων} $\mathcal B_x$. Όπου η μουσικολογική ανάλυση τεκμηριώνει ότι δύο ή περισσότερα γειτονικά ζεύγη αρχής--τέλους είναι εξίσου αποδεκτά για την ίδια εμφάνιση, λ.χ. αν ένα συνδετικό πάτημα ανήκει στην κατάληξη της προηγούμενης θέσης ή στην έναρξη της επόμενης, το μοντέλο δεν εξαναγκάζει μία μοναδική τομή, αλλά ενσωματώνει το σύνολο $\mathcal B_x$ τόσο στην εκπαίδευση όσο και στην αξιολόγηση.
\item \emph{Κατηγορία άγνωστης οικογένειας} $\varnothing_F$. Το μοντέλο λειτουργεί σε καθεστώς «ανοικτού κόσμου», διαθέτοντας ρητή κατηγορία για τμήματα που παρουσιάζουν τη δομή μουσικής θέσης αλλά ανήκουν σε οικογένεια εκτός του γνωστού λεξιλογίου εκπαίδευσης.
\item \emph{Βαθμονομημένη αποχή}. Διακριτά από την κατηγορία $\varnothing_F$, όταν τα σημειογραφικά ή τροπικά τεκμήρια ενός χωρίου είναι αμφίσημα και η μέγιστη πιθανότητα υπαγωγής πέφτει κάτω από κατώφλι που έχει βαθμονομηθεί στα δεδομένα ανάπτυξης, το σύστημα απέχει από την απόφαση αντί να παράγει αυθαίρετη ετικέτα.
\end{enumerate}

Για να εκτιμηθεί τι προσφέρει η ρητή μοντελοποίηση της θέσης πέρα από τις γενικές ιδιότητες της μελωδικής επιφάνειας και πέρα από τις σύγχρονες μεθόδους της υπολογιστικής μουσικολογίας, η Μ1 πραγματοποιεί δύο σειρές συγκρίσεων. Πρώτον, αξιολογεί μέσω δοκιμών αφαίρεσης τρεις κλασικές, ερμηνεύσιμες οικογένειες γενικών ενδείξεων μελωδικών ορίων: (α) το μοντέλο τοπικών ορίων LBDM του E.~Cambouropoulos (2001a), που εντοπίζει τοπικές ασυνέχειες στα διαστήματα και στις διάρκειες· (β) την πληροφοριακή απροσδοκησία και την εντροπία του επόμενου μελωδικού γεγονότος κατά το γνωσιακό πρότυπο IDyOM (Pearce κ.ά. 2010), που ελέγχει αν η έναρξη μιας θέσης συμπίπτει με σημείο υψηλής προβλεπτικής αβεβαιότητας και το εσωτερικό της με περιοχή σταδιακής μείωσης της εντροπίας· και (γ) την κατάτμηση βάσει μελωδικού παραλληλισμού και επανάληψης (Cambouropoulos 2006· Conklin και Anagnostopoulou 2006). Έτσι αναδεικνύεται ποια όρια θέσεων εξηγούνται ήδη από απλές σημειογραφικές ασυνέχειες ή επαναλήψεις και ποια απαιτούν ειδική γνώση της οικογένειας.

Δεύτερον, αντιπαραβάλλει το τμηματικό πλαίσιο των A/B/C με τρεις σύγχρονους συγκριτές σε κοινά έργα εκτός εκπαίδευσης: έναν αυτοπαλίνδρομο Transformer σε επίπεδο ατομικού γεγονότος, $T_{\text{evt}}$, που διαθέτει μεγάλη χωρητικότητα και μηχανισμό αυτοπροσοχής για να μαθαίνει μακρές εξαρτήσεις στην ακολουθία των νευμάτων χωρίς ρητή τμηματική βαθμολόγηση· ένα υβριδικό νευρωνικό μοντέλο κατάτμησης φράσεων Transformer--CRF (Guan κ.ά. 2018, 2025), καθώς και την παραλλαγή του με κεφαλή semi-CRF, ώστε να απομονωθεί η συμβολή της τμηματικής βαθμολόγησης από την ισχύ του νευρωνικού κωδικοποιητή· και το μη εποπτευόμενο, μη παραμετρικό μπεϋζιανό μοντέλο ιεραρχικής διαδικασίας Pitman--Yor, το οποίο εφαρμόστηκε πρόσφατα από τους V.~Lanz και J.~Hajič (2025) στην κατάτμηση του γρηγοριανού μέλους.

Η σύγκριση με το μοντέλο Pitman--Yor έχει ιδιαίτερη σημασία για την ιστορική και υπολογιστική μουσικολογία. Στη μελέτη τους για το γρηγοριανό μέλος, οι Lanz και Hajič έδειξαν ότι τα τμήματα που επάγει στατιστικά ο αλγόριθμος, βελτιστοποιώντας την απομνημόνευση και τη συμπίεση του σώματος, προβλέπουν αποτελεσματικά τον τρόπο (mode), αλλά δεν ταυτίζονται με τις παραδοσιακές φόρμουλες που η θεωρία της κεντωνοποίησης θεωρεί δομικούς λίθους του μέλους. Στη βυζαντινή ψαλτική, όπου η ύπαρξη, η ονομασία και η διδασκαλία των θέσεων μαρτυρούνται ρητά στις πηγές, η Μ1 επιτρέπει να ελεγχθεί απευθείας αυτό το ερώτημα: σε ποια σημεία οι στατιστικά επαγόμενες μονάδες συμπίπτουν με τις ιστορικά και αναλυτικά τεκμηριωμένες θέσεις του παπαδικού μέλους, και σε ποια σημεία η μουσικολογική ταυτότητα της θέσης υπερβαίνει τη στατιστική επανάληψη της επιφάνειας.

Η επιλογή της αναπαράστασης $\Rs$, η οποία θα αποτελέσει την κοινή είσοδο των επόμενων μελετών Α2 και Μ2, γίνεται αποκλειστικά στα εσωτερικά δεδομένα ανάπτυξης με προκαθορισμένο κανόνα: επιλέγεται η κωδικοποίηση του πλέγματος $U\times\Phi$ που επιτυγχάνει το υψηλότερο μέσο $F1_{\text{seg}}$ (σωστό τμήμα με σωστή οικογένεια) στις τρεις αρχιτεκτονικές A/B/C, υπό τον όρο ότι ικανοποιεί ένα ελάχιστο κατώφλι μέσης ανάκλησης οικογενειών $r_{\min}$, το οποίο ορίζεται πριν από την εκπαίδευση με βάση την πιλοτική μελέτη και ανεξάρτητες κρίσεις ειδικών. Μετά την επιλογή της $\Rs$, διεξάγεται μία φορά, χωρίς καμία ανατροφοδότηση προς την επιλογή, ο ιστορικός έλεγχος δύσκολων περιπτώσεων: η ικανότητα της επιλεγμένης αναπαράστασης να διαχωρίζει αφενός τα «δύσκολα αρνητικά» ζεύγη ομοιομετροφωνικών αλλά ιστορικά διαφορετικών θέσεων, όπως το αντικένωμα, το λύγισμα και το έτερον παρακάλεσμα στον Καραγκούνη (2003, σ.~788), και αφετέρου τα «δύσκολα θετικά» ζεύγη έντονης αναλυτικής ή τροπικής διαφοροποίησης της ίδιας θέσης (ό.π., σσ.~810--812).

%========================================================== 6
\section{Μελέτη Α2: Η μουσικολογική γνώση ως πειραματική μεταβλητή}

Αφού η Μ1 προσδιορίσει την αναπαράσταση $\Rs$, η δεύτερη μελέτη του βασικού πυρήνα θέτει ένα ερώτημα που βρίσκεται στο επίκεντρο της μουσικολογικά ενημερωμένης μηχανικής μάθησης (Musicology-Informed Machine Learning, MIML): όταν ενσωματώνουμε σε ένα υπολογιστικό μοντέλο τις διαγνωστικές σχέσεις μιας μουσικολογικής προδιαγραφής, λ.χ. απαιτούμενη βαθμίδα κατάληξης, επιτρεπτά εσωτερικά σχήματα και αποκλειόμενες τροπικές μεταβολές, η βελτίωση της αναγνώρισης οφείλεται πράγματι στο ότι το μουσικολογικό περιεχόμενο των κανόνων είναι αληθές και διαγνωστικό, ή μήπως οποιοσδήποτε δομημένος κανόνας ίδιας πολυπλοκότητας θα λειτουργούσε εξίσου καλά ως μαθηματικός περιορισμός; Και επιπλέον, τι προσθέτει η ακριβής λογική σύνθεση των κριτηρίων της προδιαγραφής πέρα από την απλή προσθετική άθροιση των επιμέρους ενδείξεων;

Στη σύγχρονη νευροσυμβολική μηχανική μάθηση, λ.χ. στα μοντέλα στενωπού εννοιών (concept bottleneck models· Koh κ.ά. 2020) και στη σημασιολογική απώλεια (semantic loss· Xu κ.ά. 2018), οι κανόνες των ειδικών θεωρούνται συνήθως δεδομένοι. Στην ιστορική μουσικολογία, ωστόσο, οι προδιαγραφές των οικογενειών αποτελούν αναθεωρήσιμες αναλυτικές υποθέσεις, ενώ πρόσφατες εργασίες έχουν δείξει ότι τα νευροσυμβολικά μοντέλα μπορούν να εκμεταλλευτούν στατιστικές «συντομεύσεις» των εννοιών (Marconato κ.ά. 2023· van Krieken κ.ά. 2024), επιτυγχάνοντας υψηλή επίδοση ακόμη και με εσφαλμένους κανόνες. Για να απαντηθεί το ζήτημα χωρίς αμφισημία, η Α2 σχεδιάζεται ως ελεγχόμενο παραγοντικό πείραμα $2\times2$, το οποίο διασταυρώνει την εγκυρότητα του περιεχομένου της γνώσης (έγκυρη προδιαγραφή $V$ έναντι ελεγχόμενα διαταραγμένης προδιαγραφής $P$) με τον τρόπο ενσωμάτωσής της στο μοντέλο (επίπεδη προσθετική χρήση $F$ έναντι δομημένης λογικής χρήσης $D$).

Για κάθε υποψήφιο τμήμα και για κάθε οικογένεια θέσεων, αυτόματοι εξαγωγείς υπολογίζουν από τη συμβολική είσοδο $\Rs$ τις στοιχειώδεις μουσικές έννοιες που χρησιμοποιούν οι προδιαγραφές (κατάληξη, προσέγγιση, εσωτερικά διαστηματικά και ρυθμικά σχήματα, ενεργές μαρτυρίες και φθορές), χωρίς πρόσβαση στην ετικέτα αναφοράς· η αξιοπιστία των εξαγωγέων ελέγχεται από τους ειδικούς ήδη στην πιλοτική μελέτη του πρώτου εξαμήνου. Πάνω στον ίδιο πίνακα εξαγόμενων εννοιών, και κρατώντας σταθερή τη βασική βαθμολογία $b_{\hat\psi}(s,f)$ του μοντέλου της Μ1 ώστε να απομονώνεται η επίδραση της παρέμβασης, συγκρίνονται δύο τρόποι σύνθεσης. Στην επίπεδη χρήση $F$, το μοντέλο συνδυάζει γραμμικά τρεις συνοπτικούς δείκτες που αντιστοιχούν στους ρόλους των γνωρισμάτων: ποσοστό ικανοποίησης αναγκαίων σχέσεων, υποστηρικτικών σχέσεων και απουσίας αποκλειστικών σχέσεων. Στη δομημένη λογική χρήση $D$, το μοντέλο αξιολογεί επιπλέον την πλήρη λογική πρόταση της προδιαγραφής, για παράδειγμα έναν κανόνα της μορφής $z_1\wedge(z_2\vee z_3)\wedge\neg z_4$, όπου απαιτείται η κατάληξη $z_1$, τουλάχιστον μία από τις δύο εναλλακτικές εσωτερικές αναπτύξεις $z_2,z_3$ και αποκλείεται η τροπική μεταβολή $z_4$, ή συνδυασμό δύο ανεξάρτητων εναλλακτικών ομάδων $(a\vee b)\wedge(c\vee d)$, ο οποίος δεν είναι γραμμικά διαχωρίσιμος από έναν απλό προσθετικό κανόνα. Οι δύο συνθήκες γράφονται
\[
\ell_{F}=b_{\hat\psi}+w^{\top}Q_f,\qquad
\ell_{D}=b_{\hat\psi}+w^{\top}Q_f-\lambda\,v_f,
\]
όπου $Q_f$ οι τρεις συνόψεις ρόλων, $w$ μη αρνητικά βάρη κοινά στις οικογένειες, και $v_f$ δείκτης τεκμηριωμένης παραβίασης της λογικής πρότασης, η οποία επιφέρει ρητή ποινή $\lambda$ στην κανονικοποίηση του semi-CRF τόσο κατά την εκπαίδευση όσο και κατά την πρόβλεψη. Επειδή στις ιστορικές πηγές ένα επιμέρους γνώρισμα μπορεί να μην είναι αξιολογήσιμο σε συγκεκριμένο χωρίο, η Α2 εφαρμόζει τριμερή λογική: η κατάσταση «άγνωστο» διακρίνεται από την «τεκμηριωμένη παραβίαση» και δεν επιφέρει ποινή στην κύρια ανάλυση.

Για να ελεγχθεί η αξία του ίδιου του μουσικολογικού περιεχομένου, οι συνθήκες $F$ και $D$ επαναλαμβάνονται αντικαθιστώντας την έγκυρη προδιαγραφή με δέκα ελεγχόμενα διαταραγμένες προδιαγραφές. Ο γεννήτορας των διαταραχών δεν δημιουργεί τυχαίο θόρυβο, αλλά κατασκευάζει ψευδο-προδιαγραφές ίδιας λογικής δομής και πολυπλοκότητας: διατηρεί το συντακτικό δένδρο του κανόνα, την τροπική εφαρμοσιμότητα, τις επικαλύψεις μεταξύ οικογενειών και τις εμπειρικές συχνότητες εμφάνισης και ελλείψεων των εννοιών, εναλλάσσοντας όμως τις μουσικές σχέσεις, λ.χ. αντικαθιστώντας την απαιτούμενη βαθμίδα κατάληξης ή το διαγνωστικό εσωτερικό σχήμα με άλλο, μουσικολογικά μη διαγνωστικό γνώρισμα ίσης συχνότητας, και πιστοποιώντας αριθμητικά ότι η διαταραγμένη προδιαγραφή δεν είναι λογικά ισοδύναμη με την αρχική. Από τον σχεδιασμό προκύπτουν δύο επιβεβαιωτικά μεγέθη πάνω στη μέση ακρίβεια κατάταξης των οικογενειών, $\mathrm{AP}_{\text{fam}}$:
\[
\Delta_{\text{know}}=\tfrac12\bigl\{S(D,V)-\bar S(D,P)+S(F,V)-\bar S(F,P)\bigr\},\qquad
\Delta_{\text{logic},V}=S(D,V)-S(F,V).
\]
Το πρώτο μετρά πόσο βελτιώνεται η αναγνώριση όταν οι κανόνες περιέχουν έγκυρη μουσικολογική γνώση αντί για ισόμορφους αλλά μουσικολογικά εσφαλμένους κανόνες· το δεύτερο μετρά τι προσθέτει η λογική σύνθεση της έγκυρης προδιαγραφής πέρα από την απλή προσθετική χρήση των ίδιων εννοιών.

Οι τέσσερις συνθήκες συγκρίνονται επιπλέον με δύο μοντέλα αναφοράς: έναν ταξινομητή $B_K$ που δέχεται τις ίδιες εξαγόμενες έννοιες και μαθαίνει ελεύθερα γραμμικά βάρη χωρίς ανθρώπινη προδιαγραφή, και έναν νευρωνικό συγκριτή $B^{\star}$ αυξημένης χωρητικότητας, που συνδυάζει εποπτευόμενη αντιπαραθετική μάθηση (SupCon· Khosla κ.ά. 2020), πρωτότυπα οικογενειών και εποπτεία εννοιών. Η μουσικολογική αξία της Α2 είναι διττή. Αν τα $\Delta_{\text{know}}$ και $\Delta_{\text{logic},V}$ είναι θετικά, υπό έλεγχο πολλαπλότητας κατά Holm, αποδεικνύεται πειραματικά ότι οι διαγνωστικές σχέσεις που εντόπισε η μουσικολογική έρευνα, όπως οι κανόνες βαθμίδας, γένους και λειτουργίας στον Gheorghiță (2009) και στον Καραγκούνη (2003), είναι αναγκαίες για τη διάκριση των θέσεων και δεν μπορούν να υποκατασταθούν από απλή στατιστική συσχέτιση. Αν σε κάποια οικογένεια ο ελεύθερος ταξινομητής $B_K$ υπερτερήσει της γραπτής προδιαγραφής, το αποτέλεσμα υποδεικνύει στον μουσικολόγο ποια διαγνωστική σχέση είχε διατυπωθεί υπερβολικά στενά ή ευρέως.

%========================================================== 7
\section{Μελέτη Μ2: Οικογενειακή ταυτότητα, διατήρηση της πραγμάτωσης και ιστορική εποπτεία από συν-εξηγήσεις}

Ενώ η Α2 εξετάζει τη γνώση που διατυπώνεται με ρητούς κανόνες, η Μ2 εξετάζει πώς μαθαίνεται η ταυτότητα μιας θέσης απευθείας από τις σχέσεις μεταξύ των εμφανίσεών της στα μουσικά έργα. Εδώ ανακύπτει ένα θεμελιώδες μουσικολογικό και υπολογιστικό δίλημμα: όταν εκπαιδεύουμε έναν νευρωνικό κωδικοποιητή να φέρνει κοντά στον χώρο αναπαράστασης δύο διαφορετικές εμφανίσεις της ίδιας οικογένειας, λ.χ. μία σύντομη και μία αναπτυγμένη εκδοχή της παρακλητικής, υπάρχει ο κίνδυνος το μοντέλο να επιτύχει την ταύτιση «σβήνοντας» κάθε εσωτερική μελωδική και ρυθμική διαφορά τους. Για τον μουσικολόγο, όμως, μια αναπαράσταση είναι χρήσιμη μόνο όταν απαντά ταυτόχρονα σε δύο ερωτήματα: ποια θέση είναι αυτή και πώς ακριβώς πραγματώθηκε στο συγκεκριμένο χωρίο.

Για να επιλύσει αυτό το δίλημμα, ο τμηματικός κωδικοποιητής της Μ2 δέχεται την αναπαράσταση $\Rs$, υπολογίζει μια τοπική κατάσταση $h_t$ για κάθε μελωδικό γεγονός εντός των συμφραζομένων του και εξάγει για κάθε υποψήφιο τμήμα $[s,e]$ δύο διακριτές αναπαραστάσεις: την προβολή οικογενειακής ταυτότητας $z_{\text{id}}$, όπου εφαρμόζεται η βαθμολόγηση του semi-CRF και η έλξη των συγγενών εμφανίσεων, και την αναπαράσταση πραγμάτωσης $z_{\text{real}}$, μαζί με την ακολουθία των τοπικών καταστάσεων $h_t$, η οποία προστατεύεται μέσω παράλληλου αυτοεποπτευόμενου στόχου ανακατασκευής αποκρυμμένων συμβολικών γεγονότων, τμημάτων ή συλλαβών (κατά το πρότυπο του MusicBERT· Zeng κ.ά. 2021).

Ένα εξαιρετικό ιστορικό παράδειγμα για τον έλεγχο αυτού του διαχωρισμού ανάμεσα στη δομική ταυτότητα και στην τροπικά προσδιορισμένη πραγμάτωση προσφέρει η ανάλυση του οκτάηχου Κοινωνικού του Πέτρου Μπερεκέτη από τον Gheorghiță (2009, §7.3, σσ.~230--237, παραδείγματα 11--12, με βάση τον κώδικα Ξηροποτάμου 380, φφ.~402v--403r). Το παλαιό σημειογραφικό πρότυπο είναι γραμμένο χωρίς καμία μαρτυρία, ώστε να μπορεί να ψαλεί στους οκτώ ήχους, ενώ στην αναλυτική εξήγηση του Γρηγορίου Πρωτοψάλτη αποδίδεται σε έξι διακριτές μελωδικές εκδοχές· ο ίδιος ο Γρηγόριος εξηγεί, για το αντίστοιχο Χερουβικό και με τον ίδιο τρόπο για το Κοινωνικό, ότι ο πρώτος και ο πλάγιος του πρώτου, όπως και ο δεύτερος και ο πλάγιος του δευτέρου, «οδεύουσιν ομοίαν κλίμακα και καταλήξεις». Ο Gheorghiță χωρίζει το πρώτο μέρος του Κοινωνικού σε δεκαπέντε δομικές ενότητες και διαπιστώνει ότι ο σκελετός και οι μελωδικοί πυλώνες της παλαιάς γραφής διατηρούνται πλήρως στη Νέα Μέθοδο, ενώ αναπραγματώνονται μελωδικά ανάλογα με το τροπικό περιβάλλον κάθε ήχου. Από υπολογιστική σκοπιά, το έργο λειτουργεί ως σπάνιο φυσικό πείραμα, όπου ο μελοποιός, το παλαιό πρότυπο και ο εξηγητής παραμένουν σταθεροί και μεταβάλλεται κατά βάση η τροπική παράμετρος. Εντάσσοντάς το ως στοχευμένο διαγνωστικό σύνολο ελέγχου στη Μ2, μπορούμε να μετρήσουμε αν ο τμηματικός κωδικοποιητής αναγνωρίζει την κοινή ταυτότητα των δεκαπέντε δομικών ενοτήτων στις έξι εξηγήσεις μέσω του $z_{\text{id}}$, διατηρώντας ταυτόχρονα τις τροπικές διαφοροποιήσεις κάθε ήχου στον κλάδο πραγμάτωσης $z_{\text{real}}$.

Στο πρώτο σκέλος της μελέτης (Μ2a), το οποίο στηρίζεται αποκλειστικά στο υλικό της Νέας Μεθόδου και ανήκει στον βασικό πυρήνα, συγκρίνονται τρεις κλιμακωτές συνθήκες. Η συνθήκη B0 είναι το συμβολικό semi-CRF της Μ1 χωρίς εκπαιδευμένο νευρωνικό κωδικοποιητή και δείχνει τι επιτυγχάνει από μόνη της η άμεση περιγραφή της γραφής. Η συνθήκη B1 προσθέτει τον προεκπαιδευμένο κωδικοποιητή με εποπτεία οικογενειακής ετικέτας και στόχο ανακατασκευής, χωρίς σχεσιακό όρο ζευγών. Η συνθήκη B2 διατηρεί την ίδια αρχιτεκτονική, αρχικοποίηση και προϋπολογισμό με τη B1, προσθέτοντας έναν ελκτικό όρο θετικών ζευγών που φέρνει κοντά στον χώρο $z_{\text{id}}$ ανεξάρτητα επισημειωμένες εμφανίσεις της ίδιας οικογένειας από διαφορετικά έργα:
\[
\mathcal L_{\text{B2}}=\mathcal L_{\text{seg}}+\lambda_f\,\mathcal L_{\text{fam}}+\lambda_r\,\mathcal L_{\text{rec}}+\lambda_p\,\mathcal L_{\text{pair}},\qquad \lambda_p=0\ \text{στη B1}.
\]
Η αξιολόγηση της Μ2a σε έργα εκτός εκπαίδευσης είναι διπλή: μετρά αφενός το κέρδος στην οικογενειακή αναγνώριση, $\Delta_{\text{famrel}}=\mathrm{AP}_{\text{fam}}(\text{B2})-\mathrm{AP}_{\text{fam}}(\text{B1})$, και αφετέρου τον δείκτη διατήρησης της πραγμάτωσης $Q_{\text{real}}$. Ο δείκτης αυτός ελέγχει, με τρεις προκαθορισμένες δοκιμές ανίχνευσης, αν η αναπαράσταση της B2 εξακολουθεί να διατηρεί προσβάσιμη την πληροφορία της συγκεκριμένης εμφάνισης: ανακατασκευή κρυμμένων σημείων, διάκριση σύντομων και αναπτυγμένων πραγματώσεων της ίδιας οικογένειας και πρόβλεψη του μήκους και της συλλαβικής κατανομής, με προκαθορισμένο όριο μη κατωτερότητας έναντι της B1. Επιπλέον, επειδή δύο εμφανίσεις της ίδιας οικογένειας σε Χερουβικά ή Κοινωνικά ενδέχεται να βρίσκονται πάνω στην ίδια λέξη του λειτουργικού κειμένου, λ.χ. «τὰ Χερουβίμ» ή «Αἰνεῖτε», η Μ2a επαναλαμβάνει τη σύγκριση αφαιρώντας το κειμενικό στρώμα $L$ από τον κλάδο ταυτότητας ($\Rs_{\setminus L}$), ώστε να αποκλειστεί το ενδεχόμενο το μοντέλο να αναγνωρίζει τη λέξη αντί για τη μουσική θέση (Sturm 2014).

Στο δεύτερο, ιστορικό σκέλος της μελέτης (Μ2b: σύγκριση B2--EISE), εξετάζεται τι μεταβάλλεται στη μαθημένη οργάνωση της παραλλαγής όταν τα θετικά ζεύγη εκπαίδευσης δεν είναι απλώς εμφανίσεις της ίδιας οικογένειας από διαφορετικές συνθέσεις, όπου συνυπάρχουν διαφορές μελοποιού, κειμένου και τροπικού τόπου, αλλά ιστορικά επαληθευμένες συν-εξηγήσεις του ίδιου παλαιότερου χωρίου, λ.χ. η απόδοση του Γρηγορίου και η απόδοση του Χουρμουζίου για το ίδιο χωρίο ενός Κοινωνικού του Δανιήλ. Στη συνθήκη EISE (Exegesis-Invariant Segmental Encoder), ο κωδικοποιητής εκπαιδεύεται με την ίδια συνάρτηση απώλειας της B2, αλλά τα θετικά ζεύγη προέρχονται από τις συν-εξηγήσεις, ενώ τα ζεύγη της B2 αντιστοιχίζονται ως προς τον αριθμό των ακμών, την έκθεση κάθε εμφάνισης, τον βαθμό του γράφου, την επιφανειακή απόσταση στο στρώμα $N$ και τον υπολογιστικό προϋπολογισμό.

Για να κριθεί ποια από τις δύο πηγές εποπτείας, η κοινή αναλυτική κατηγορία στη B2 ή η κοινή ιστορική αναφορά στην EISE, οργανώνει με μουσικά εγκυρότερο τρόπο την εσωτερική παραλλαγή κάθε θέσης, η Μ2b χρησιμοποιεί ένα ανεξάρτητο κριτήριο από το επίπεδο της παρατήρησης, το οποίο κανένα από τα δύο μοντέλα δεν συναντά κατά την εκπαίδευση: τη χρονική μεταβλητότητα σύγχρονων ψαλτών, $\Gamma_f(u)$, όταν ψάλλουν τα ίδια χωρία (ελάχιστο σκέλος Μ4-min). Για κάθε οικογένεια $f$, οι εμφανίσεις ευθυγραμμίζονται σε κοινή εσωτερική μουσική συντεταγμένη $u$ (αρχή, εσωτερικές ζώνες, κατάληξη) χωρίς χρήση εκτελεστικού χρόνου. Από τις τοπικές καταστάσεις $h(u)$ κάθε μοντέλου $m\in\{\text{B0},\text{B1},\text{B2},\text{EISE}\}$, μετά από λεύκανση της συνδιακύμανσης που εκτιμάται αποκλειστικά στο σύνολο εκπαίδευσης, εξάγεται το μαθημένο ιστορικό προφίλ παραλλαγής
\[
\hat\tau_{f,m,\text{hist}}(u)=\sqrt{\operatorname{tr}\widehat\Sigma_{f,m}(u)/d_h}.
\]
Παράλληλα, από τις ηχογραφήσεις πολλών ανεξάρτητων ψαλτών στα ίδια χωρία υπολογίζεται ανά περιοχή $u$ η διασπορά του κανονικοποιημένου λογαριθμικού χρονισμού $\Gamma_f(u)$, αφού αφαιρεθεί η συνολική αγωγή της ηχογράφησης και η τοπική πυκνότητα των σημειογραφικών αξιών. Το κύριο συγκριτικό μέγεθος είναι
\[
\Delta_{\text{hist}}=S_{\text{EISE}}-S_{\text{B2}},
\]
όπου $S_m$ η μέση συσχέτιση κατάταξης Spearman μεταξύ $\hat\tau_{f,m,\text{hist}}(u)$ και $\Gamma_f(u)$ ανά οικογένεια. Το μέγεθος αυτό ελέγχει μια γόνιμη μουσικολογική υπόθεση: αν τα εσωτερικά σημεία μιας θέσης όπου δύο εξηγητές επέλεξαν να αναπτύξουν, να συμπτύξουν ή να διαφοροποιήσουν την αναλυτική γραφή συμπίπτουν συστηματικά με τις ζώνες όπου η σύγχρονη ψαλτική πράξη αφήνει μεγαλύτερη χρονική ελευθερία στον εκτελεστή (πρβλ. Cambouropoulos 2001a· Gingras κ.ά. 2016). Η σύγκριση δεν υποθέτει ότι οι σημερινοί ψάλτες αναπαράγουν αναλλοίωτο το ιστορικό άκουσμα των αρχών του 19ου αιώνα· αντιμετωπίζει τη γραπτή εξηγητική διαφορά και τη σύγχρονη εκτελεστική ελευθερία ως δύο ανεξάρτητες μαρτυρίες για το ποια μέρη μιας θέσης αποτελούν τον σταθερό δομικό της πυρήνα και ποια επιδέχονται ελαστικότητα. Όπως ορίστηκε στην §3, η πιλοτική εκπαίδευση της EISE προχωρά μόλις η απογραφή τεκμηριώσει επιλέξιμα ιστορικά ζεύγη, ενώ ο επιβεβαιωτικός στατιστικός έλεγχος του $\Delta_{\text{hist}}$ ενεργοποιείται μόνο εφόσον πληρούνται ταυτόχρονα τα προκαθορισμένα κριτήρια επάρκειας ζευγών ($G_{\text{pair}}$), ηχογραφήσεων ($G_{\text{perf}}$) και βαθμονομημένης στατιστικής ακρίβειας ($G_{\text{precision}}$).

%========================================================== 8
\section{Ανεξάρτητη μουσικολογική επικύρωση (Μ3a), έλεγχος σχολικής μεροληψίας και εφαρμογή σε πλήρη έργα (Μ5-min)}

Όλες οι υπολογιστικές μελέτες του πυρήνα θα κινδύνευαν από κυκλικό συλλογισμό αν οι προδιαγραφές των οικογενειών και οι επισημειώσεις των χωρίων προέρχονταν αποκλειστικά από τον ερευνητή που σχεδιάζει τα μοντέλα. Ο κίνδυνος είναι ακόμη πιο συγκεκριμένος στη βυζαντινή μουσικολογία: ο υποψήφιος είναι ενεργός ψάλτης με μακρόχρονη μαθητεία στη διδακτική γραμμή του Σίμωνος Καρά και του Λυκούργου Αγγελόπουλου. Η πρακτική αυτή εξοικείωση είναι πολύτιμη για τη σύνταξη των αρχικών υποθέσεων και την αναγνώριση της φρασεολογίας των θέσεων, θα μπορούσε όμως να εισαγάγει συστηματική σχολική απόκλιση, ιδίως σε ζητήματα όπου διαφορετικές ψαλτικές ή ακαδημαϊκές παραδόσεις αναλύουν διαφορετικά τα όρια ή τη λειτουργία των ποιοτικών χαρακτήρων. Για να εξασφαλιστεί η αντικειμενικότητα της έρευνας, η Μ3a εισάγει πρωτόκολλο ανεξάρτητης μουσικολογικής επικύρωσης και ελέγχου κυκλικότητας με πέντε κανόνες:
\begin{enumerate}
\item \emph{Αποκλεισμός των κρίσεων του υποψηφίου από την αξιολόγηση.} Οι επισημειώσεις και οι αναλύσεις του υποψηφίου καταχωρίζονται στο μητρώο αποκλειστικά ως υλικό ορισμού και ανάπτυξης και δεν προσμετρώνται ποτέ στα σύνολα τελικής αξιολόγησης ή επιβεβαίωσης.
\item \emph{Τελική έκδοση των προδιαγραφών πριν από τον έλεγχο.} Κάθε προδιαγραφή αναπτύσσεται και αναθεωρείται μόνο στο υλικό ανάπτυξης· πριν παραδοθεί στους ανεξάρτητους κριτές και πριν εκπαιδευτούν τα τελικά μοντέλα (μήνας 12), καταχωρίζεται σε τελική, αριθμημένη έκδοση, η οποία δεν μεταβάλλεται στη συνέχεια. Αν μια οικογένεια αποτύχει στον έλεγχο συμφωνίας των ειδικών, δεν διορθώνεται εκ των υστέρων με βάση τις προβλέψεις των μοντέλων, αλλά καταγράφεται ως αμφίσημη ή υποψήφια κατηγορία.
\item \emph{Τυφλή κρίση από ειδικούς διαφορετικής προέλευσης.} Συμμετέχουν τουλάχιστον δύο ανεξάρτητοι ειδικοί της βυζαντινής μουσικολογίας και ψαλτικής τέχνης, οι οποίοι δεν συνέταξαν τις προδιαγραφές και από τους οποίους τουλάχιστον ο ένας προέρχεται από διδακτική ή ερευνητική παράδοση εκτός της σχολής του υποψηφίου. Σε πρώτη φάση, οι ειδικοί κρίνουν τυφλά, χωρίς να βλέπουν ούτε τις προδιαγραφές ούτε τις εξόδους των αλγορίθμων, τα όρια, την ένταξη νέων χωρίων (ένταξη, μη ένταξη, ανεπίλυτο) και δοκιμασίες τριάδων («ποιο από τα τρία χωρία δεν ανήκει με τα άλλα δύο;»), όπου το ξένο χωρίο έχει επιλεγεί ώστε να παρουσιάζει υψηλή επιφανειακή ή μετροφωνική ομοιότητα. Σε δεύτερη φάση, εφαρμόζουν τις προδιαγραφές στην τελική τους έκδοση, ελέγχοντας ένα προς ένα τα διαγνωστικά τους γνωρίσματα.
\item \emph{Η διαφωνία των ειδικών ως μουσικολογικό εύρημα.} Ακολουθώντας τη μεθοδολογία της υπολογιστικής γλωσσολογίας για την ανθρώπινη διακύμανση στην επισημείωση (Artstein και Poesio 2008· Plank 2022), οι διαφωνίες δεν εξαλείφονται σιωπηρά μέσω ψηφοφορίας πλειοψηφίας. Κάθε διαφωνία διατηρείται στο μητρώο και ταξινομείται κατά την αιτία της: ασάφεια ορίων, διαφορά αναλυτικής κλίμακας μέρους--όλου, διαφορετική τροπική ανάγνωση, π.χ. σε ενδιάμεσες μαρτυρίες ἀπὸ παραλλαγῆς ή ἀπὸ μέλους (πρβλ. Troelsgård 2011, σ.~69· Alexandru και Troelsgård 2020, σσ.~178--179), διαφωνία ως προς το διαγνωστικό γνώρισμα, ανεπάρκεια πηγής ή διάσταση ιστορικής ονομασίας. Με τον τρόπο αυτό χαρτογραφούνται τα σημεία όπου η ίδια η ψαλτική παράδοση επιτρέπει πολλαπλές αναλυτικές αναγνώσεις.
\item \emph{Ομαδοποίηση έργων και πηγών.} Κατά τον διαχωρισμό των δεδομένων σε εκπαίδευση, ανάπτυξη και έλεγχο, ολόκληρα έργα, χειρόγραφα ή έντυπα της ίδιας γενεαλογικής γραμμής και όλες οι συν-εξηγήσεις του ίδιου χωρίου τοποθετούνται στην ίδια ομάδα, ώστε κανένα τμήμα του ίδιου Χερουβικού ή της ίδιας εξηγητικής αλυσίδας να μη βρίσκεται ταυτόχρονα στην εκπαίδευση και στην αξιολόγηση.
\end{enumerate}

Ως καταληκτικό βήμα του βασικού πυρήνα, η Μ5-min μεταφέρει τα επικυρωμένα εργαλεία από το επίπεδο της μεμονωμένης θέσης στην ανάλυση ολόκληρων έργων εκτός εκπαίδευσης. Για κάθε πλήρες Χερουβικό ή Κοινωνικό παράγεται ένας αναλυτικός χάρτης έργου, ο οποίος αποτυπώνει στη συνεχή ροή του μέλους τις αναγνωρισμένες εμφανίσεις θέσεων, με τα όρια, την οικογένεια και τη βαθμονομημένη βεβαιότητα της αναγνώρισης, διατηρώντας ορατές τόσο τις αναγνωρισμένες συνδετικές θέσεις (κατά τους πίνακες 8--9 του Gheorghiță 2010) όσο και το μη αποδοθέν συνδετικό υλικό ή τις ανεπίλυτες περιοχές. Η ελάχιστη περιγραφική σύνοψη κάθε έργου ποσοτικοποιεί δύο μουσικολογικές διαστάσεις: (α) την επιλογή οικογενειών $D_{\text{fam}}$, δηλαδή το προφίλ του λεξιλογίου θέσεων που χρησιμοποιεί η σύνθεση, και (β) την ενδοοικογενειακή πραγμάτωση $D_{\text{within}}$, η οποία μετρά, πάνω στις τοπικές καταστάσεις του κωδικοποιητή και σε ένα ερμηνεύσιμο χρονικό διαγνωστικό $(Q_x,p_x,d_{\text{time}})$, πώς αναπτύσσει ή συμπτύσσει ο μελοποιός τις θέσεις κάθε οικογένειας. Κάθε ποσοτικός δείκτης διατηρεί αμφίδρομη σύνδεση με τα μουσικά χωρία της παρτιτούρας, ώστε ο μουσικολόγος να μπορεί να ελέγξει αμέσως από ποιες φράσεις προκύπτει οποιαδήποτε στατιστική διαφοροποίηση.

\subsection*{Ενιαίο παράδειγμα: το καταληκτικό κύλισμα από την πηγή έως τον χάρτη του έργου}

Για να φανεί πώς συνεργάζονται οι μελέτες του πυρήνα πάνω σε ένα συγκεκριμένο μουσικολογικό αντικείμενο, ας παρακολουθήσουμε την πορεία της θέσης με κύλισμα. Στο ΠΕ0 (μήνες 1--6), το κύλισμα επιλέγεται, μαζί με την παρακλητική και το πίασμα, ως μία από τις τρεις πιλοτικές οικογένειες με βάση τις αναλύσεις του Gheorghiță (2009), επειδή έχει σαφή καταληκτική λειτουργία και συχνή εμφάνιση στα Κοινωνικά. Η πρώτη έκδοση της προδιαγραφής του δηλώνει ως αναγκαία την κατάληξη σε καταληκτική βαθμίδα του ενεργού ήχου, ως υποστηρικτικά τα χαρακτηριστικά νευματικά σχήματα που προηγούνται της κατάληξης, και ως αποκλειστικό μια τροπική μεταβολή που μετατοπίζει τη βάση· δύο ειδικοί επισημειώνουν ανεξάρτητα τις εμφανίσεις του σε μικρό σύνολο έργων. Στη Μ1, το μοντέλο εντοπίζει στη συνεχή γραφή τα σημεία όπου αρχίζει και τελειώνει η θέση και την αναγνωρίζει ως κύλισμα, ή απέχει όπου τα τεκμήρια δεν επαρκούν· η σύγκριση με τη στατιστική κατάτμηση δείχνει αν οι καταλήξεις με κύλισμα αναδύονται και ως στατιστικά επαναλαμβανόμενα τμήματα. Στην Α2, η έγκυρη προδιαγραφή συγκρίνεται με αλλοιωμένες εκδοχές όπου η απαιτούμενη κατάληξη τοποθετείται σε λάθος βαθμίδα· αν η αναγνώριση υποχωρεί, η βαθμίδα κατάληξης αποδεικνύεται διαγνωστική. Στη Μ2a, οι εμφανίσεις του κυλίσματος σε διαφορετικά έργα διδάσκουν την ταυτότητά του, ενώ ο δείκτης $Q_{\text{real}}$ ελέγχει ότι η αναπαράσταση εξακολουθεί να διακρίνει τη σύντομη από την αναπτυγμένη πραγμάτωση και τη διατονική από τη χρωματική. Στη Μ3a, ειδικοί κρίνουν τυφλά νέα χωρία, και μια τριάδα με επιφανειακά όμοια μη καταληκτική θέση ελέγχει αν η κρίση στηρίζεται στη λειτουργία. Τέλος, στη Μ5-min, ο χάρτης ενός πλήρους Κοινωνικού δείχνει πού χρησιμοποιείται το κύλισμα και πόσο διαφοροποιούνται οι πραγματώσεις του από ήχο σε ήχο. Αν οι ιστορικές προϋποθέσεις ικανοποιηθούν, το ίδιο κύλισμα μπορεί να εξεταστεί και στο αντίστροφο πρόβλημα της §9.1: αν από την εξήγηση ανακτάται ότι ο παλαιός μάρτυρας έγραφε κύλισμα και όχι άλλο σημάδι με την ίδια μετροφωνία.

%========================================================== 9
\section{Ιστορικές, παλαιογραφικές, γνωσιακές και συγκριτικές προεκτάσεις υπό προϋποθέσεις}

Πάνω στην υποδομή του βασικού πυρήνα, η πρόταση έχει σχεδιάσει μια σειρά αυτοτελών ερευνητικών επεκτάσεων, καθεμία από τις οποίες ενεργοποιείται με ρητό κριτήριο επάρκειας των αντίστοιχων πηγών, χωρίς να επιβαρύνει ή να θέτει σε κίνδυνο την ολοκλήρωση του κορμού.

\subsection{Το αντίστροφο παλαιογραφικό πρόβλημα (ATL-INV) και οι κώδικες εξηγήσεων θέσεων}

Η σχέση ανάμεσα στην παλαιότερη, συνοπτικότερη σημειογραφία και στην αναλυτική απόδοση της Νέας Μεθόδου υπήρξε επί δεκαετίες ένα από τα πλέον συζητημένα ζητήματα της βυζαντινής μουσικολογίας. Όπως συνοψίζει ο Ε. Γιαννόπουλος (2014, σσ.~160--161· πρβλ. Στάθης 1978· Καράς 1982· Lingas 2003· Αλεξάνδρου 2010· Αρβανίτης 2010, 2022), το κεντρικό ερώτημα αφορά το αν η εξήγηση αποτέλεσε πιστή αναλυτική καταγραφή, με φωνητικά σημάδια, του μέλους που διδασκόταν προφορικά και υποδηλωνόταν από τις μεγάλες υποστάσεις, ή αν εισήγαγε νέους μελωδικούς σχηματισμούς. Η διατριβή δεν ισχυρίζεται ότι ένα τόσο σύνθετο ιστορικό ζήτημα μπορεί να λυθεί με έναν υπολογιστικό δείκτη· δείχνει όμως ότι ένα συγκεκριμένο, κρίσιμο μέρος του μπορεί να τεθεί ως ποσοτικά ελέγξιμο αντίστροφο πρόβλημα.

Ενώ η ιστορική διαδικασία της εξήγησης κινείται από το παλαιό σημάδι προς την αναλυτική καταγραφή, το αντίστροφο ερώτημα αντιστρέφει την κατεύθυνση: αν διαθέτουμε μόνο την αναλυτική εξήγηση της Νέας Μεθόδου και γνωρίζουμε το τροπικό και προσωδιακό περιβάλλον του χωρίου, με ποια πιθανότητα μπορούμε να ανακτήσουμε ποιο σημάδι ή σύμπλεγμα σημαδιών είχε καταγράψει ο παλαιός χειρόγραφος μάρτυρας; Η δυσκολία του ερωτήματος έγκειται στο ότι διαφορετικές γραπτές επιλογές της παλαιότερης σημειογραφίας μοιράζονται την ίδια μετροφωνία: ως απλό παράδειγμα, το ολίγον και η πεταστή μπορούν να δηλώνουν την ίδια διαστηματική άνοδο κατά μία φωνή, ενώ οι άφωνες μεγάλες υποστάσεις, όπως το αντικένωμα, το λύγισμα, το πίασμα ή το κύλισμα, συνδυάζονται με κοινή ποσοτική βάση χωρίς να μεταβάλλουν από μόνες τους τη διαστηματική διαδρομή (Alexandru 2012· Καραγκούνης 2003, σ.~788).

Ο πειραματικός σχεδιασμός εκμεταλλεύεται ακριβώς αυτή την ιδιότητα. Συμβολίζουμε με $H$ την παλαιά ετικέτα που διαβάζεται στον συγκεκριμένο μάρτυρα, με $M$ τη μετροφωνική της κλάση, με $Z$ το τροπικό και προσωδιακό περιβάλλον και με $X_G,X_C$ τις εξηγήσεις του ίδιου χωρίου, λ.χ. του Γρηγορίου και του Χουρμουζίου. Ως σύνολο υποψηφίων $\mathcal H(M,Z)$ ορίζονται όλες οι ιστορικά επιλέξιμες σημειογραφικές επιλογές με την ίδια μετροφωνία στο ίδιο περιβάλλον, ενώ η αλήθεια αναφοράς προέρχεται από ανεξάρτητη, τυφλή παλαιογραφική ανάγνωση του μάρτυρα, χωρίς πρόσβαση στην εξήγηση. Ένα εμπρόσθιο ιστορικό κανάλι $q_\phi(X_G\mid H,M,Z)$ και η εκ των προτέρων κατανομή $p(H\mid M,Z)$ δίνουν την αντίστροφη κατανομή
\[
p_\phi(h\mid x_G,m,z)=\frac{q_\phi(x_G\mid h,m,z)\,p(h\mid m,z)}{\sum_{h'\in\mathcal H(m,z)}q_\phi(x_G\mid h',m,z)\,p(h'\mid m,z)}.
\]
Επειδή όλοι οι υποψήφιοι μιας μετροφωνικής κλάσης δίνουν εξ ορισμού την ίδια διαστηματική ανάγνωση, όση πληροφορία για το παλαιό σημάδι παρέχουν η μετροφωνία και το περιβάλλον την αποδίδει ήδη το μοντέλο αναφοράς. Κάθε πρόσθετο προβλεπτικό κέρδος από τη γνώση της εξήγησης οφείλεται επομένως στη μελική οργάνωση της αναλυτικής γραφής: στον τρόπο με τον οποίο αναπτύσσεται το μέλισμα, στους χρονικούς και ποιοτικούς χαρακτήρες της Νέας Μεθόδου και στην κατανομή τους στη συλλαβή.

Πάνω στον ίδιο στόχο πραγματοποιούνται τρεις διαδοχικές συγκρίσεις, σε bits πληροφορίας ανά ιστορική ετικέτα και σε έργα εκτός εκπαίδευσης:
\[
\begin{aligned}
\Delta_{\text{rec}}&=\mathrm{CE}(H\mid M,Z)-\mathrm{CE}(H\mid M,Z,X_G),\\
\Delta_{\text{rec}}^{\text{rule}}&=\mathrm{CE}\bigl(H\mid M,Z,\widehat H_{\text{rule}}(X_G)\bigr)-\mathrm{CE}(H\mid M,Z,X_G),\\
\Delta_{\text{2ex}}&=\mathrm{CE}(H\mid M,Z,X_G)-\mathrm{CE}(H\mid M,Z,X_G,X_C).
\end{aligned}
\]
Η $\Delta_{\text{rec}}$ μετρά πόσο βελτιώνεται η ανάκτηση του παλαιού σημαδιού όταν γνωρίζουμε την πρώτη εξήγηση σε σχέση με μόνο τη μετροφωνία και το περιβάλλον· η $\Delta_{\text{rec}}^{\text{rule}}$ μετρά αν το πιθανοτικό μοντέλο αντλεί από την εξήγηση περισσότερη διακριτική πληροφορία από όση κωδικοποιούν ήδη οι δημοσιευμένοι κανόνες εξήγησης, οι οποίοι, από τις πραγματείες του Αποστόλου Κώνστα, του Ψάχου και του Στάθη, διατυπώνονται εκ των προτέρων ως ισχυρός συγκριτής· και η $\Delta_{\text{2ex}}$ μετρά πόση πρόσθετη πληροφορία προσφέρει μια δεύτερη εξήγηση του ίδιου χωρίου, όπου σώζονται παράλληλες αποδόσεις Γρηγορίου και Χουρμουζίου. Όταν για το ίδιο χωρίο υπάρχουν περισσότεροι του ενός παλαιοί μάρτυρες, η συμφωνία τους ορίζει εμπειρικά το ανώτατο όριο σταθερότητας του ιστορικού στόχου. Έτσι η μουσικολογία αποκτά έναν ποσοτικό «χάρτη ανακτησιμότητας», ο οποίος δείχνει ποιες διακρίσεις της παλαιάς γραφής, και ποιες ενέργειες ποιοτικών σημαδιών και μεγάλων υποστάσεων, αφήνουν συστηματικά ανιχνεύσιμο αποτύπωμα στην αναλυτική γραφή και ποιες συγχωνεύονται.

Στο σημείο αυτό αναδεικνύεται και η μεθοδολογική σημασία του τέταρτου παράγοντα μεταβολής που εντοπίζει ο Gheorghiță (2009, σσ.~250--251), δηλαδή του ερυθρού ή μελανού χρώματος του χειρονομικού σημαδιού στους παλαιούς κώδικες, το οποίο κατά τον Απόστολο Κώνστα καθορίζει αν η θέση θα αναπτυχθεί απλή ή σύνθετη, συντομότερη ή μακρύτερη. Επειδή στη Νέα Μέθοδο το παλαιό χειρονομικό σημάδι και το χρώμα του δεν μεταφέρονται ως γραφικά σύμβολα, η πρόταση τηρεί κανόνα καταγραφής ανά επίπεδο τεκμηρίωσης. Στις μελέτες που τροφοδοτούνται αποκλειστικά από τη Νέα Μέθοδο (Μ1, Μ2a, Μ5-min), το χρώμα του παλαιού σημαδιού δεν εισάγεται ως χαρακτηριστικό εισόδου· η πιθανή επίδρασή του μπορεί να αποτυπώνεται μόνο έμμεσα, στη συνολική σημειογραφική χρονική ανάπτυξη $Q_x$ και στην εσωτερική κατανομή του χρόνου $p_x$ της αναλυτικής πραγμάτωσης, και αυτό αντιμετωπίζεται ως υπόθεση προς έλεγχο. Όταν όμως το χωρίο συνδέεται με συγκεκριμένο παλαιό μάρτυρα στις ιστορικές επεκτάσεις, η παρουσία ερυθρού ή μελανού σημαδιού καταγράφεται στο επίπεδο της πηγής, και μόνο όπου η ανάγνωση είναι βέβαιη. Αυτό επιτρέπει να εξεταστεί ποσοτικά αν η διαφορά χρώματος στον παλαιό κώδικα προβλέπει τη διαφορά έκτασης στην εξήγηση, στο εμπρόσθιο κανάλι, και αντίστροφα αν από τη μορφολογία της αναλυτικής ανάπτυξης μπορεί να ανακτηθεί πιθανοτικά όχι μόνο η ταυτότητα αλλά και η απλή ή σύνθετη λειτουργία της υποκείμενης μεγάλης υπόστασης.

Όσον αφορά τις πηγές αυτής της επέκτασης, πέρα από τους ονομαστικά προσδιορισμένους κώδικες της παράλληλης μελέτης των κυριακών Κοινωνικών στον Gheorghiță (2009, σσ.~128--129), ιδιαίτερα προνομιακό υλικό προσφέρουν οι χειρόγραφες συλλογές εξηγήσεων σε επίπεδο μεμονωμένης θέσης. Όπως έδειξε ο Γιαννόπουλος (2014, σσ.~171--181), ο κώδικας 35 της Μονής Αγίας Αναστασίας Φαρμακολυτρίας (περίπου 1815--1819) είναι οργανωμένος σε ζεύγη με τις ενδείξεις «π» για τη θέση στην παλαιότερη γραφή και «ε» για την αναλυτική της απόδοση, παραθέτοντας συχνά δύο εξηγήσεις της ίδιας θέσης, μία σε αναλυτικότερη παλαιά γραφή και μία στη Νέα Μέθοδο· ο Γιαννόπουλος εξετάζει αν ο γραφέας αντιπαραβάλλει στην πραγματικότητα τρία ιστορικά στρώματα: τη γραφή του Πέτρου Πελοποννησίου, τις πρωτοεξηγήσεις του Γρηγορίου και του Χουρμουζίου και τη Νέα Μέθοδο. Στο ίδιο είδος πηγών ανήκει και η συλλογή του Ακακίου Μεγασπηλαιώτη του έτους 1800 (ΜΜΑ «Λίλιαν Βουδούρη», χφ.~1). Εάν οι παπαδικές μαρτυρίες τέτοιων συλλογών αποδειχθούν αριθμητικά περιορισμένες κατά την απογραφή, το μοντέλο ανάκτησης μπορεί να εκπαιδευτεί και να βαθμονομηθεί αρχικά στο πλουσιότερο στιχηραρικό τμήμα των ίδιων κωδίκων και κατόπιν να εφαρμοστεί στα παπαδικά χωρία ως ελεγχόμενη μελέτη περίπτωσης.

\subsection{Μορφολογία των εξηγητικών μετασχηματισμών (ATL) και πιθανοτική μοντελοποίηση της πραγμάτωσης ($\Mth$ και ECSI)}

Δύο συγγενείς επεκτάσεις συμπληρώνουν τη μελέτη της παραλλαγής. Η πρώτη, η μάθηση τεκμηριωμένων μετασχηματισμών (Attested Transformation Learning, ATL), εξετάζει απευθείας τη σχέση μεταξύ δύο παράλληλων εξηγήσεων του ίδιου χωρίου. Χρησιμοποιώντας έναν στοχαστικό μετατροπέα ακολουθιών που επεκτείνει τη μουσική απόσταση επεξεργασίας με πράξεις κατάτμησης, όπου μία κίνηση αναλύεται σε περισσότερες, και συνένωσης, όπου περισσότερες κινήσεις συμπτύσσονται (Mongeau και Sankoff 1990· Ristad και Yianilos 1998· Giraud κ.ά. 2014), η ATL στοιχίζει τις δύο εξηγήσεις μέσα σε κάθε συλλαβή, αθροίζοντας πάνω σε όλες τις επιτρεπτές στοιχίσεις $A$:
\[
q_\theta(X_C\mid X_G,F,Z)=\sum_{A\in\mathcal A(X_G,X_C)}q_\theta(A,X_C\mid X_G,F,Z),
\]
και εκτιμά αν οι κανονικότητες ανάπτυξης και σύμπτυξης που μαθαίνονται σε ορισμένες οικογένειες μεταφέρονται προβλεπτικά σε οικογένειες που είχαν αποκρυφτεί από την εκπαίδευση ($\Delta_{\text{ATL}}$). Μια συγκεκριμένη μουσικολογική υπόθεση προς έλεγχο είναι η διαπίστωση του Gheorghiță για τον Μπερεκέτη, ότι η εξήγηση διατηρεί τον σκελετό του παλαιού προτύπου και τον «διαστέλλει»· η ATL μπορεί να δείξει σε ποιες εσωτερικές περιοχές της θέσης συγκεντρώνεται η διαστολή. Ακόμη και ένα αρνητικό αποτέλεσμα στη δοκιμή μεταφοράς έχει άμεση μουσικολογική σημασία, διότι τεκμηριώνει ότι οι εξηγητικοί μετασχηματισμοί είναι εξειδικευμένοι ανά οικογένεια θέσεων ή ανά εξηγητή.

Η δεύτερη επέκταση, ο πιθανοτικός τελεστής πραγμάτωσης $\Mth$, εξετάζει τη σχέση τύπου και πραγμάτωσης μέσα στη Νέα Μέθοδο: όταν γνωρίζουμε την οικογένεια $F$ μιας θέσης και το τοπικό, τροπικό και προσωδιακό της περιβάλλον $C$, πόσο προβλέψιμη γίνεται η ακριβής μελωδική της μορφή $X$ σε σύγκριση με ένα μοντέλο που γνωρίζει μόνο το περιβάλλον;
\[
\Delta_{\text{family}}=\mathrm{CE}(X\mid C)-\mathrm{CE}(X\mid F,C).
\]
Το μοντέλο επιτρέπει να διακριθεί ποσοτικά το μουσικολογικά επιτρεπτό, ό,τι ικανοποιεί την προδιαγραφή της οικογένειας, από το εμπειρικά τυπικό ή προτιμώμενο στο εξεταζόμενο σώμα: μια σπάνια αλλά νόμιμη πραγμάτωση εντοπίζεται μέσω του δείκτη τυπικότητας $-\log_2 q_\theta(x\mid F,C)$, προσφέροντας στον μουσικολόγο εργαλείο για τον εντοπισμό ιδιαίτερων μελοποιητικών επιλογών. Σε παράλληλη αρχιτεκτονική επέκταση της Α2 (ECSI: Evidence-Conditioned Selective Invariance), η λογική προδιαγραφή κάθε οικογένειας μεταγλωττίζεται σε ακριβές πιθανοτικό κύκλωμα (Darwiche και Marquis 2002· Xu κ.ά. 2018), ώστε να ελέγχεται η υπόθεση ανεξαρτησίας των διαγνωστικών εννοιών (van Krieken κ.ά. 2024).

\subsection{Πρόσληψη στην ακρόαση (Μ3b, Μ3c), πλήρης εκτελεστική ανάλυση (Μ4) και η μελοποιητική πρακτική του Νικολάου Πρωτοψάλτη Σμύρνης (πλήρης Μ5)}

Τρεις επιπλέον κατευθύνσεις συνδέουν τη συμβολική ανάλυση με τον ακροατή, τον εκτελεστή και την ιστορική προσωπογραφία ενός μελοποιού.

\emph{Γνωσιακή επικύρωση στην ακρόαση (Μ3c και Μ3b).} Εξετάζεται αν οι οικογένειες θέσεων που ορίστηκαν μουσικολογικά λειτουργούν ως ενεργές γνωσιακές κατηγορίες κατά την ακρόαση, συνεχίζοντας το ερώτημα του Mavromatis (2005b· 2019). Στην πρώτη συνεδρία (Μ3c: ιεραρχική μουσική προσδοκία), ακροατές διαφορετικής εμπειρίας, έμπειροι ψάλτες και μουσικοί χωρίς βυζαντινή εκπαίδευση, ακούν αποσπάσματα και κρίνουν πόσο απρόσμενη είναι η συνέχειά τους. Μέσω πιθανοτικού μοντέλου μεταβλητού συμφραζομένου (Prediction by Partial Matching, PPM, στο πνεύμα του IDyOM· Pearce 2018), ελέγχεται αν το πληροφοριακό περιεχόμενο της επόμενης οικογένειας θέσεων $I_F$ εξηγεί τις ανθρώπινες κρίσεις απροσδοκησίας πέρα από ό,τι εξηγούν τα μεμονωμένα μελωδικά γεγονότα $I_{\text{event}}$, με ισορροπημένα αντιδιαγώνια ερεθίσματα (υψηλό $I_F$ με χαμηλό $I_{\text{event}}$ και αντίστροφα). Στη δεύτερη συνεδρία (Μ3b: αντιληπτική κατηγοριοποίηση), οι ίδιοι ακροατές κρίνουν τριάδες αποδόσεων επιλέγοντας εκείνη που διαφέρει περισσότερο, με ελεγχόμενη την επιφανειακή ακουστική και μελωδική απόσταση, ώστε να διαπιστωθεί αν οργανώνουν τις πραγματώσεις βάσει της οικογενειακής τους ταυτότητας (πρβλ. Cambouropoulos 2001b· Volk και van Kranenburg 2012).

\emph{Πλήρης εκτελεστική μελέτη (Μ4).} Πέρα από το ελάχιστο κριτήριο χρονισμού της Μ2b, η πλήρης Μ4 μελετά συστηματικά τη σχέση γραπτής θέσης και φωνητικής εκτέλεσης σε υπάρχουσες ηχογραφήσεις, όπως του σώματος ΔΑΜΑΣΚΗΝΟΣ, και σε νέες ηχογραφήσεις που μπορούν να γίνουν στις εργαστηριακές υποδομές ακουστικής του Τμήματος: πώς κατανέμεται η διασπορά του χρονισμού, η καταληκτική επιμήκυνση και η μικροδιαστηματική διαμόρφωση του τονικού ύψους ($F_0$) στις εσωτερικές ζώνες κάθε οικογένειας, και κατά πόσον οι εκτελεστικές αυτές επιλογές διαφοροποιούνται ανά διδακτική γραμμή του ψάλτη. Οι μετρήσεις αυτές μένουν ανεξάρτητο επίπεδο τεκμηρίωσης· δεν υποκαθιστούν τη μουσικολογική κρίση ούτε αποδεικνύουν ιστορική συνέχεια.

\emph{Πλήρης Μ5 και εφαρμογή στον Νικόλαο Πρωτοψάλτη Σμύρνης.} Όταν το σώμα αναφοράς $\Omega$ διαθέτει επαρκή κάλυψη συνεχών μεταβάσεων και ανεξάρτητων μελοποιών, η Μ5-min επεκτείνεται με δύο επιπλέον διαστάσεις: τη διαδοχή των οικογενειών $D_{\text{seq}}$, μέσω του μέσου πληροφοριακού περιεχομένου $I_\Omega$ των αλυσίδων θέσεων και της διάκρισης μεταξύ αυτούσιας επανάληψης αλυσίδας και ανασυνδυασμού γνωστών θέσεων (κατά τον Καραγκούνη 2005, σσ.~175--176), και τη λειτουργική τοποθέτηση $D_{\text{func}}$, δηλαδή την κατανομή των οικογενειών ανά δομικό, τροπικό και κειμενικό τόπο. Το πλήρες αυτό τετραδιάστατο προφίλ $(D_{\text{fam}},D_{\text{within}},D_{\text{seq}},D_{\text{func}})$ εφαρμόζεται στη μελοποιία του Νικολάου Πρωτοψάλτη Σμύρνης του 19ου αιώνα (Τόμα 2018), έναντι χρονολογικά ελεγχόμενου σώματος αναφοράς, διακρίνοντας το προγενέστερο στρώμα του 18ου αιώνα, που σώζεται μέσω εξηγήσεων, από το σύγχρονο στρώμα του 19ου αιώνα και ελέγχοντας την επίδραση κοινού εξηγητή. Έτσι εξετάζεται αν η ιδιοτυπία του Νικολάου έγκειται στην επιλογή διαφορετικών θέσεων, στον ασυνήθιστο σειριακό συνδυασμό τους ή στον εσωτερικό τρόπο ανάπτυξής τους, ενώ η δημοσιευμένη μορφολογική ανάλυση του Χ.~Ιωαννίδη (2022) για το Κοινωνικό του Νικολάου «Εἰς πᾶσαν τὴν γῆν» σε ήχο πλ. δ΄, σε αντιπαραβολή με τον Πέτρο Λαμπαδάριο και τον Χουρμούζιο, λειτουργεί ως προκαθορισμένο, ανεξάρτητο σημείο εξωτερικού ελέγχου.

\subsection{Διαπαραδοσιακή συγκριτική επέκταση: η βυζαντινή θέση και οι ποπέβκι του ζναμένιου μέλους}

Ως απαιτητική δοκιμή της γενικευσιμότητας του μεθοδολογικού πλαισίου πέρα από την ελληνική ψαλτική, η πρόταση προβλέπει συγκριτική πιλοτική εφαρμογή στο ζναμένιο μέλος της ρωσικής ορθόδοξης παράδοσης, όπως σώζεται στα ιστορικά χειρόγραφα αλλά και στη ζωντανή πρακτική κοινοτήτων Παλαιοπίστων, όπως στην περιοχή της λίμνης Πείψι (Grigorjeva 2023· Parfentiev και Parfentieva 2023· Gusev κ.ά. 2020). Η νευματική γραφή του ζναμένιου μέλους κατάγεται ιστορικά από τη βυζαντινή και το ρεπερτόριό του οργανώνεται μέσα στο σύστημα της οκταηχίας με βάση επανεμφανιζόμενους μελωδικούς τύπους, τις ποπέβκι, ενώ οι εκτενείς μελισματικοί σχηματισμοί που γράφονταν συνοπτικά ως φίτες αποδίδονται στις πηγές και αναλυτικά, ως ραζβόντι, προσφέροντας μια εντυπωσιακή δομική αναλογία προς τη σχέση μεγάλης υπόστασης και εξήγησης. Χωρίς να ταυτίζει ιστορικά ή μουσικολογικά την ποπέβκα με τη βυζαντινή θέση, η επέκταση εξετάζει αν η ίδια υπολογιστική αρχιτεκτονική, ο πιθανοτικός πυρήνας στοίχισης και η δομημένη σύνθεση ECSI-D, μπορεί να αναγνωρίσει την ταυτότητα μέσα στην παραλλαγή όταν κάθε παράδοση ορίζει τις δικές της ανεξάρτητες μουσικολογικές προδιαγραφές. Οι έως τώρα έλεγχοι χρησιμοποιούν συνθετικά δεδομένα· η εφικτότητα της σύγκρισης εξαρτάται από την πρόσβαση στις πηγές, την επισημείωση και τη συνεργασία με ειδικούς της ιστορικής σημειογραφίας.

%========================================================== 10
\section{Συνοπτικός απολογισμός των μουσικολογικών παραδοτέων και όρια της έρευνας}

Ανεξάρτητα από την έκβαση των επιμέρους στατιστικών υποθέσεων, η ολοκλήρωση της διατριβής θα αποδώσει στην έρευνα της βυζαντινής μουσικολογίας ένα σύνολο συγκεκριμένων, ανοικτά τεκμηριωμένων εργαλείων και πορισμάτων:
\begin{enumerate}
\item Ένα μητρώο οικογενειών θέσεων στα παπαδικά Χερουβικά και Κοινωνικά της Νέας Μεθόδου, με αριθμημένες εκδόσεις, όπου κάθε οικογένεια συνοδεύεται από ρητή προδιαγραφή (αναγκαίες, υποστηρικτικές και αποκλειστικές σχέσεις, εναλλακτικές διαμορφώσεις, κανόνες τροπικού και προσωδιακού περιβάλλοντος) και πλήρη βιβλιογραφική και πηγολογική τεκμηρίωση.
\item Ένα ψηφιακά κωδικοποιημένο και επισημειωμένο σώμα παπαδικών μελών, ιχνηλάσιμο ως τις πηγές και συμβατό με την αναπτυσσόμενη προσαρμογή του προτύπου MEI για τα βυζαντινά νεύματα (Alexandru κ.ά. 2024· Thomae κ.ά. 2026), στο οποίο καταγράφονται όχι μόνο οι εμφανίσεις θέσεων, αλλά και τα εναλλακτικά αποδεκτά όρια και οι ταξινομημένες κατά αιτία διαφωνίες των ανεξάρτητων ειδικών.
\item Μια πειραματική αποτίμηση των ίδιων των μουσικολογικών κριτηρίων, μέσω των μελετών Μ1, Α2 και Μ2a, η οποία δείχνει σε ποια κλίμακα ανάγνωσης της γραφής αναγνωρίζονται καλύτερα οι θέσεις, πού η στατιστική κατάτμηση αποκλίνει από την παραδοσιακή θέση, ποια διαγνωστικά γνωρίσματα είναι πράγματι αναγκαία για τη διάκριση ομοιομετροφωνικών σχηματισμών και πώς κατανέμεται η εσωτερική παραλλαγή κάθε οικογένειας.
\item Αναλυτικούς χάρτες πλήρων έργων (Μ5-min), οι οποίοι επιτρέπουν την επαναλήψιμη επισκόπηση της μορφολογικής διάρθρωσης Χερουβικών και Κοινωνικών, και, εφόσον πληρωθούν τα κριτήρια επάρκειας των ιστορικών και εκτελεστικών επεκτάσεων, μια απογραφή των σωζόμενων συν-εξηγήσεων, έναν χάρτη παλαιογραφικής ανακτησιμότητας των παλαιών σημαδιών από τις εξηγήσεις τους, μια ποσοτική μορφολογία της εξηγητικής πράξης του Γρηγορίου και του Χουρμουζίου και μια συγκριτική αποτίμηση της μελοποιίας του Νικολάου Πρωτοψάλτη Σμύρνης.
\end{enumerate}

Κλείνοντας, είναι αναγκαίο να υπογραμμιστούν τα επιστημολογικά όρια της πρότασης. Η υπολογιστική ανάλυση δεν υποκαθιστά την παλαιογραφική αυτοψία, τη φιλολογική κριτική των κωδίκων ή την ερμηνευτική κρίση του μουσικολόγου· ούτε επιδιώκει να αποφανθεί ποια σύγχρονη ψαλτική σχολή ή θεωρητική προσέγγιση της εξήγησης είναι η ορθή. Αντίθετα, θέτει την υπολογιστική μοντελοποίηση στην υπηρεσία της μουσικολογικής τεκμηρίωσης: υποχρεώνει κάθε αναλυτική υπόθεση να διατυπωθεί ρητά, μετρά με προκαθορισμένους κανόνες πόσο μακριά φθάνουν οι διαγνωστικοί κανόνες σε νέα έργα και αναδεικνύει με ακρίβεια τα σημεία όπου η ζωντανή μουσική παράδοση υπερβαίνει την τυποποίηση.

%========================================================== BIB
\section*{Βιβλιογραφικές αναφορές}
\small
\setlength{\parindent}{0pt}
\subsection*{Ελληνόγλωσση βιβλιογραφία}
\begin{list}{}{\leftmargin=1.5em\itemindent=-1.5em\itemsep=1pt}
\item Αλεξάνδρου, Μ. (2010). \emph{Εξηγήσεις και μεταγραφές της βυζαντινής μουσικής: Σύντομη εισαγωγή στον προβληματισμό τους}. Θεσσαλονίκη: University Studio Press.
\item Αλεξάνδρου, Μ. (2017). \emph{Παλαιογραφία Βυζαντινής Μουσικής}. Αθήνα: Σύνδεσμος Ελληνικών Ακαδημαϊκών Βιβλιοθηκών (Κάλλιπος). Ιδίως το κεφάλαιο «Μεσοβυζαντινή σημειογραφία (Δ΄): Χειρονομία, Μεγάλα σημάδια, Θέσεις».
\item Αναστασίου, Γ. Γ. (2005). \emph{Τα Κρατήματα στην Ψαλτική Τέχνη} (Μελέται 12). Αθήνα: Ίδρυμα Βυζαντινής Μουσικολογίας.
\item Αποστολόπουλος, Θ. Κ. (2002). \emph{Ο Απόστολος Κώνστας ο Χίος και η συμβολή του στη θεωρία της μουσικής τέχνης} (Μελέται 4). Αθήνα: Ίδρυμα Βυζαντινής Μουσικολογίας.
\item Αρβανίτης, Ι. (2010). \emph{Ο ρυθμός των εκκλησιαστικών μελών μέσα από την παλαιογραφική έρευνα και την εξήγηση της παλαιάς σημειογραφίας} (διδακτορική διατριβή, 2 τόμοι). Κέρκυρα: Ιόνιο Πανεπιστήμιο. \url{https://doi.org/10.12681/eadd/21885}
\item Γιαννόπουλος, Ε. Σ. (2014). Η εξήγηση των παλαιών θέσεων της ψαλτικής. Νέες μαρτυρίες. \emph{Ανατολής το Περιήχημα} 1, 159 κ.ε.
\item Δεσπότης, Σ. Κ. (2007). Ερμηνευτικές προσεγγίσεις στο μουσικό υλικό της ελληνικής ψαλτικής τέχνης. \emph{Γρηγόριος ο Παλαμάς} 818, 417--428.
\item Ιωαννίδης, Χ. (2022). Η σύνθεση του Νικολάου Πρωτοψάλτου Σμύρνης στο Κοινωνικό «Εἰς πᾶσαν τὴν γῆν» σε ήχο πλ. δ΄ και ο βαθμός εξάρτησής της από την προγενέστερη παράδοση της εκκλησιαστικής μελοποιίας. \emph{Θεολογία} 93(4), 173--193.
\item Καραγκούνης, Κ. Χ. (2003). \emph{Η Παράδοση και Εξήγηση του Μέλους των Χερουβικών της Βυζαντινής και μετα-Βυζαντινής Μελοποιίας} (Μελέται 7). Αθήνα: Ίδρυμα Βυζαντινής Μουσικολογίας.
\item Καραγκούνης, Κ. Χ. (2005). \emph{Παραλειπόμενα περὶ τοῦ Χερουβικοῦ Ὕμνου}. Βόλος: Γραφικές Τέχνες Αλέξανδρος Ξουράφας.
\item Καράς, Σ. Ι. (1982). \emph{Μέθοδος της Ελληνικής Μουσικής: Θεωρητικόν} (τ. Α΄--Β΄). Αθήνα: Σύλλογος προς Διάδοσιν της Εθνικής Μουσικής.
\item Σπανουδάκης, Δ. Π. (2017). \emph{Σύγχρονα μοντέλα ανάλυσης Βυζαντινής Μουσικής: Το φαινόμενο της μελισματικότητας στα κεκραγάρια} (διδακτορική διατριβή, 3 τόμοι). Θεσσαλονίκη: ΑΠΘ. \url{https://doi.org/10.12681/eadd/42193}
\item Στάθης, Γ. Θ. (1978). \emph{Η εξήγησις της Παλαιάς Βυζαντινής Σημειογραφίας} (Μελέται 2). Αθήνα: Ίδρυμα Βυζαντινής Μουσικολογίας (γ΄ έκδ. 1993).
\item Τόμα, Ι. Γ. (2018). \emph{Νικόλαος Γεωργίου, πρωτοψάλτης Σμύρνης: Βίος, ψαλτικό έργο, μουσικολογικές προσεγγίσεις} (μεταπτυχιακή διπλωματική εργασία). Θεσσαλονίκη: ΑΠΘ.
\item Χρύσανθος ὁ ἐκ Μαδύτων (1832). \emph{Θεωρητικὸν Μέγα τῆς Μουσικῆς}. Τεργέστη: Michele Weis.
\item Ψάχος, Κ. Α. (1917/1978). \emph{Η Παρασημαντική της Βυζαντινής Μουσικής} (β΄ επαυξημένη έκδ., επιμ. Γ. Χατζηθεοδώρου). Αθήνα: Διόνυσος.
\end{list}
\subsection*{Ξενόγλωσση βιβλιογραφία}
\begin{list}{}{\leftmargin=1.5em\itemindent=-1.5em\itemsep=1pt}
\item Alexandru, M. (2012). Some thoughts on the so-called great hypostases of the Middle Byzantine notation. \emph{Musicology Papers} 27(1), 47--58.
\item Alexandru, M. \& Troelsgård, C. (2020). The “Elements of the Papadikê” and modality features in Byzantine chant. \emph{Series Musicologica Balcanica} 1(1), 169--204.
\item Alexandru, M., De Luca, E. \& Fujinaga, I. (2024). Brainstorming on the encoding of Eastern neumes: Middle-Byzantine notation as a case-study. \emph{Series Musicologica Balcanica} 5, 177--205. \url{https://doi.org/10.26262/smb.v0i5.10414}
\item Artstein, R. \& Poesio, M. (2008). Inter-coder agreement for computational linguistics. \emph{Computational Linguistics} 34(4), 555--596.
\item Arvanitis, I. (2022). The transcriptions of the Three Teachers: An old or a new sound? \emph{Epistēmēs Metron Logos} 5, 143--149. \url{https://doi.org/10.12681/eml.28005}
\item Cambouropoulos, E. (2001a). The Local Boundary Detection Model (LBDM) and its application in the study of expressive timing. \emph{Proceedings of ICMC 2001}.
\item Cambouropoulos, E. (2001b). Melodic cue abstraction, similarity, and category formation: A formal model. \emph{Music Perception} 18(3), 347--370.
\item Cambouropoulos, E. (2006). Musical parallelism and melodic segmentation: A computational approach. \emph{Music Perception} 23(3), 249--268.
\item Conklin, D. \& Anagnostopoulou, C. (2006). Segmental pattern discovery in music. \emph{INFORMS Journal on Computing} 18(3), 285--293.
\item Conomos, D. E. (επιμ.) (1985). \emph{The Treatise of Manuel Chrysaphes, the Lampadarios} (MMB, Corpus Scriptorum de Re Musica 2). Wien: Verlag der Österreichischen Akademie der Wissenschaften.
\item Darwiche, A. \& Marquis, P. (2002). A knowledge compilation map. \emph{Journal of Artificial Intelligence Research} 17, 229--264.
\item Delviniotis, D. \& Kouroupetroglou, G. (2011). DAMASKINOS: The prototype corpus of Greek Orthodox ecclesiastical chant voices. \emph{Proceedings of Crossroads: Greece as an Intercultural Pole of Musical Thought and Creativity}, Θεσσαλονίκη.
\item Gheorghiță, N. (2009). \emph{Chinonicul duminical în perioada post-bizantină (1453--1821). Liturgică și muzică}. București: Sophia.
\item Gheorghiță, N. (2010). Observations on the technique of transcription (exegesis) into the New Method of analytical music notation of the Sunday koinonikon of the 18th century. \emph{Studii și Cercetări de Istoria Artei. Teatru, Muzică, Cinematografie} 4(48), 9--26.
\item Gingras, B., Pearce, M. T., Goodchild, M., Dean, R. T., Wiggins, G. A. \& McAdams, S. (2016). Linking melodic expectation to expressive performance timing and perceived musical tension. \emph{Journal of Experimental Psychology: Human Perception and Performance} 42(4), 594--609.
\item Giraud, M., Déguernel, K. \& Cambouropoulos, E. (2014). Fragmentations with pitch, rhythm and parallelism constraints for variation matching. Στο \emph{Sound, Music, and Motion: CMMR 2013} (LNCS 8905), 298--312. Springer.
\item Grigorjeva, V. (2023). Principles of stratification of melodic formulas of Old Russian out-of-ichos chants. \emph{St.~Tikhon's University Review, Series V: Christian Art} 50, 9--23.
\item Guan, X., Dong, Z., Liu, H. \& Li, Q. (2025). Improving phrase segmentation in symbolic folk music: A hybrid model with local context and global structure awareness. \emph{Entropy} 27(5), 460.
\item Guan, Y., Zhao, J., Qiu, Y., Zhang, Z. \& Xia, G. (2018). Melodic phrase segmentation by deep neural networks. arXiv:1811.05688.
\item Gusev, V. D., Bakhmutova, I. V., Miroshnichenko, L. A. \& Titkova, T. N. (2020). Possible approaches to deciphering Russian ancient znamenny chant. \emph{Journal of New Music Research} 49(3), 298--306.
\item Hannick, C. \& Wolfram, G. (επιμ.) (1985). \emph{Gabriel Hieromonachos: Abhandlung über den Kirchengesang} (MMB, Corpus Scriptorum de Re Musica 1). Wien: Verlag der Österreichischen Akademie der Wissenschaften.
\item Khosla, P. κ.ά. (2020). Supervised contrastive learning. \emph{Advances in Neural Information Processing Systems} 33, 18661--18673.
\item Koh, P. W. κ.ά. (2020). Concept bottleneck models. \emph{Proceedings of ICML 2020}, PMLR 119, 5338--5348.
\item Kong, L., Dyer, C. \& Smith, N. A. (2016). Segmental recurrent neural networks. \emph{Proceedings of ICLR 2016}.
\item Lanz, V. \& Hajič jr., J. (2025). Gregorian melody, modality, and memory: Segmenting chant with Bayesian nonparametrics. \emph{Proceedings of ISMIR 2025}, 638--646.
\item Lingas, A. (2003). Performance practice and the politics of transcribing Byzantine chant. \emph{Acta Musicae Byzantinae} 6, 56--76.
\item Marconato, E., Teso, S., Vergari, A. \& Passerini, A. (2023). Not all neuro-symbolic concepts are created equal: Analysis and mitigation of reasoning shortcuts. \emph{Advances in Neural Information Processing Systems} 36.
\item Mavromatis, P. (2005b). \emph{The Echoi of Modern Greek Church Chant in Written and Oral Transmission: A Computational Model and Its Cognitive Implications} (διδακτορική διατριβή). Eastman School of Music.
\item Mavromatis, P. (2019). Hierarchical structural patterns in improvised music: Implications for cognition. \emph{Society for Music Theory, 42nd Annual Meeting, Program and Abstracts}, 56--57.
\item Mongeau, M. \& Sankoff, D. (1990). Comparison of musical sequences. \emph{Computers and the Humanities} 24(3), 161--175.
\item Neanes Project (χ.χ.). \emph{Neanes: A Byzantine chant scorewriter} [λογισμικό ανοικτού κώδικα]. \url{https://github.com/neanes/neanes}
\item Parfentiev, N. P. \& Parfentieva, N. V. (2023). Stroganovs' master of church-singing art Faddey Subotin at the court of Tsar Alexei Mikhailovich. \emph{Journal of Siberian Federal University. Humanities \& Social Sciences} 16(7), 1112--1126.
\item Pearce, M. T. (2018). Statistical learning and probabilistic prediction in music cognition: Mechanisms of stylistic enculturation. \emph{Annals of the New York Academy of Sciences} 1423(1), 378--395.
\item Pearce, M. T., Müllensiefen, D. \& Wiggins, G. A. (2010). The role of expectation and probabilistic learning in auditory boundary perception: A model comparison. \emph{Perception} 39(10), 1365--1391.
\item Plank, B. (2022). The “problem” of human label variation: On ground truth in data, modeling and evaluation. \emph{Proceedings of EMNLP 2022}, 10671--10682.
\item Ristad, E. S. \& Yianilos, P. N. (1998). Learning string-edit distance. \emph{IEEE Transactions on Pattern Analysis and Machine Intelligence} 20(5), 522--532.
\item Sarawagi, S. \& Cohen, W. W. (2004). Semi-Markov conditional random fields for information extraction. \emph{Advances in Neural Information Processing Systems} 17, 1185--1192.
\item Siklafidis, N. \& Alexandru, M. (2021). Font design of Psaltic (Byzantine) notation for Greek musical repertoires. \emph{Music Encoding Conference Proceedings 2021}.
\item St.~Anthony's Greek Orthodox Monastery (χ.χ.). \emph{Papadic Formulae}. \url{https://stanthonysmonastery.org/pages/papadic-formulae}
\item Sturm, B. L. (2014). A simple method to determine if a music information retrieval system is a “horse”. \emph{IEEE Transactions on Multimedia} 16(6), 1636--1644.
\item Thomae, M., Alexandru, M., De Luca, E., Aslanidi, M., Bazmadelis, A., Fujinaga, I., Ioannidou, A. \& Siklafidis, N. (2026). A long-awaited desideratum: MEI encoding of Byzantine neumes (Chrysanthine notation). IAML Congress, Θεσσαλονίκη.
\item Troelsgård, C. (2011). \emph{Byzantine Neumes: A New Introduction to the Middle Byzantine Musical Notation} (MMB, Subsidia 9). Copenhagen: Museum Tusculanum Press.
\item Tsuboi, Y., Kashima, H., Mori, S., Oda, H. \& Matsumoto, Y. (2008). Training conditional random fields using incomplete annotations. \emph{Proceedings of COLING 2008}, 897--904.
\item van Krieken, E., Minervini, P., Ponti, E. M. \& Vergari, A. (2024). On the independence assumption in neurosymbolic learning. \emph{Proceedings of ICML 2024}, PMLR 235, 49078--49097.
\item Volk, A. \& van Kranenburg, P. (2012). Melodic similarity among folk songs: An annotation study on similarity-based categorization in music. \emph{Musicae Scientiae} 16(3), 317--339.
\item Xu, J., Zhang, Z., Friedman, T., Liang, Y. \& Van den Broeck, G. (2018). A semantic loss function for deep learning with symbolic knowledge. \emph{Proceedings of ICML 2018}, PMLR 80, 5502--5511.
\item Zeng, M., Tan, X., Wang, R., Ju, Z., Qin, T. \& Liu, T.-Y. (2021). MusicBERT: Symbolic music understanding with large-scale pre-training. \emph{Findings of ACL-IJCNLP 2021}, 791--800.
\end{list}

\end{document}
